% concatenated from main.tex, _definitions.tex, abstract.tex, Introduction.tex, problem_formulation.tex, methodology.tex, results.tex
%%%%%%%%%%%%%%%%%%%%%%%%%%%%%%%%%%%%%%%%%%%%%%%%%%%%%%%%%%%%%%%%%%%%%%%%%%%%%%%%
%2345678901234567890123456789012345678901234567890123456789012345678901234567890
%        1         2         3         4         5         6         7         8
%        1         2         3         4         5         6         7         8

\documentclass[letterpaper, 10 pt, conference]{ieeeconf}  % Comment this line out
                                                          % if you need a4paper
%\documentclass[a4paper, 10pt, conference]{ieeeconf}      % Use this line for a4
                                                          % paper

\IEEEoverridecommandlockouts                              % This command is only
                                                          % needed if you want to
                                                          % use the \thanks command
\overrideIEEEmargins
% \addtolength{\topmargin}{5pt}
\usepackage{times} % assumes new font selection scheme installed
\usepackage{amsmath,amsfonts} % assumes amsmath package installed
\usepackage{amssymb}  % assumes amsmath package installed
\usepackage{tabularx}
\usepackage[normalem]{ulem}
\usepackage{dsfont}
\usepackage{graphicx}
\usepackage{subcaption}
\usepackage[dvipsnames,usenames]{xcolor}
\usepackage[colorlinks,%
linkcolor=BrickRed,%
filecolor=BrickRed,%
citecolor=RoyalPurple,%
]{hyperref}
% \usepackage{color}
% \usepackage{todonotes}
\usepackage{algorithm}
\usepackage{algorithmic}
% \usepackage{subfigure}
% \usepackage{subcaption}
\usepackage{float}
% \usepackage{thmtools}
% \usepackage{mathtools}
\newcommand{\bh}[1]{{\color{red} #1}}
\newcommand{\at}[1]{{\color{blue} #1}}
\newcommand{\mj}[1]{{\color{green} #1}}
\DeclareFontEncoding{LS1}{}{}
\DeclareFontSubstitution{LS1}{stix}{m}{n}
\DeclareSymbolFont{stixletters}{LS1}{stix}{m}{it}
\DeclareMathAccent{\cev}{\mathord}{stixletters}{"91}

\newcommand{\calF}{{\mathcal F}}
\newcommand{\calT}{{\mathcal T}}

\def\tilW{\widetilde W}
\def\tilZ{\widetilde Z}
\def\hap{\hat{p}}
\def\haU{\widehat{U}}
\def\ddx{\frac{\ud}{\ud x}}
\def\ddt{\frac{\ud}{\ud t}}
\def\Re{\mathbb{R}}
\def\R{\mathbb{R}}
\def\T{[0,2\pi]}
\def\KL{\,\text{\rm KL}}
\def\argmin{\mathop{\text{\rm arg\,min}}}
\def\ind{\text{\rm\large 1}}
\def\varble{\,\cdot\,}
%This is p_n: \def\pxzn{p_{\text{\tiny$ X| Z^n$}}}
\def\pxz{p_{\text{\tiny$ X| Z$}}}
\def\pzx{p_{\text{\tiny$ Z | X$}}}
\def\pdzx{p_{\text{\tiny $ \ud Z | X$}}}
\def\pz{p_{\text{\tiny$Z$}}}
\def\pdz{p_{\text{\tiny$\ud Z$}}}
\def\ps{p_{\text{\tiny$S$}}}
\def\tilh{\tilde h}
\def\tilz{\tilde z}
\def\tilp{\tilde p}
\def\Hil{\mathbf H}
\def\Ltwo{\mathbf{L}^2}
\newcommand{\Real}{\mathbb R}
\def\ddx{\frac{\ud}{\ud x}}
\def\ddt{\frac{\ud}{\ud t}}
\def\frakp{\mathfrak{p}}
\def\frakh{\mathfrak{h}}
\def\frakz{\mathfrak{z}}
\def\solp{\frakp}
\def\solh{\frakh}
\def\solz{\frakz}
\def\varble{\,\cdot\,}
\def\bfInov{\bfmath{I}}
\def\Inov{I} %So we can change symbol easily later
\def\pop{\hbox{\tiny (POP)}}
\def\Lemma#1{Lemma~\ref{#1}}
% \def\problem#1{Problem~\ref{#1}}
\def\Proposition#1{Prop.~\ref{#1}}
\def\Theorem#1{Thm.~\ref{#1}}
\def\Corollary#1{Cor.~\ref{#1}}
\def\Sec#1{Sec.~\ref{#1}}
\def\Fig#1{Fig.~\ref{#1}}
\def\Appendix#1{Appendix.~\ref{#1}}
\def\notes#1{\marginpar{\tiny #1}\typeout{Notes!
Notes!
Notes!
}}
\renewcommand{\notes}[1]{\typeout{notes!}}
\def\FRAC#1#2#3{\genfrac{}{}{}{#1}{#2}{#3}}
\def\half{{\mathchoice{\FRAC{1}{1}{2}}%
{\FRAC{2}{1}{2}}%
{\FRAC{3}{1}{2}}%
{\FRAC{4}{1}{2}}}}
\def\prrr{{\cal T}}
\def\gen{{\cal D}}
\def\bbbone{{\mathchoice {\rm 1\mskip-4mu l} {\rm 1\mskip-4mu l}
{\rm 1\mskip-4.5mu l} {\rm 1\mskip-5mu l}}}
\def\ind{\bbbone}
\newcommand{\field}[1]{\mathbb{#1}}
\newcommand{\tr}{\mbox{tr}}
\def\posRe{\field{R}_+}
\def\Re{\field{R}}
\def\k{{\sf K}}
\def\Co{\field{C}}
\def\Sec#1{Sec.~\ref{#1}}
%\def \eqdef {\mathrel {\mathop {=\joinrel=}\limits^{def}}}
\def\eqdef{\mathrel{:=}}
\def\defeq{\mathrel{=:}}
\def\transpose{{\hbox{\rm\tiny T}}}
\def\clB{{\cal B}}
\def\clE{{\cal E}}
\def\clL{{\cal L}}
\def\clP{{\cal P}}
\def\clZ{{\cal Z}}
\def\clZd{{\cal Z}^{{}\, \hbox{\tiny $\blacksquare$}}}
\def\clW{{\cal W}}
\def\Sec#1{Sec~\ref{#1}}
\def\bfmath#1{{\mathchoice{\mbox{\boldmath$#1$}}%
{\mbox{\boldmath$#1$}}%
{\mbox{\boldmath$\scriptstyle#1$}}%
{\mbox{\boldmath$\scriptscriptstyle#1$}}}}
\def\bfInov{\bfmath{I}}
\def\Inov{I} %So we can change symbol easily later
\def\bfmX{\bfmath{X}}
\def\bfmZ{\bfmath{Z}}
\def\bfmU{\bfmath{U}}
\def\bfxi{\bfmath{xi}}


\newcommand{\img}{\text{Im}}
\newcommand{\barC}{\bar{C}}
\newcommand{\barc}{\bar{c}}
% \def\E{\text{E}}
\def\E{{\sf E}}
\def\Expect{{\mathbb E}}
\newcommand{\avgT}{\lim_{T\to\infty}\frac{1}{T}\!\!\!\int_0^T}
\newcommand{\avgTt}{\lim_{T\to\infty}\frac{1}{T}\!\!\!\int_t^T}
\def\intsw{\int_{\Omega}\!\int_0^{2\pi}}
\newcommand{\avgNni}{\frac{1}{N}\sum_{j\neq i}^N}
\newcommand{\avgN}{\frac{1}{N}\sum_{j=1}^N}
\newcommand{\Oscr}{\mathcal O}
\newcommand{\Lscr}{\mathcal L}
\newcommand{\dtheta}{\theta\theta}
\newcommand{\ptheta}{\partial_{\theta}}
\newcommand{\pthetaB}{\partial^2_{\dtheta}}
\newcommand{\pt}{\partial_{t}}
\newcommand{\zero}{\mathbf{0}}
\def\clC{{\cal C}}
\def\clX{{\cal X}}
\def\Prob{{\sf P}}
\def\Probsub{{\sf P\! }}

\def\state{{\sf X}}
\def\spm#1{\notes{\textcolor{blue}{SPM: #1}}}
\def\pgm#1{\notes{\textcolor{blue}{PGM: #1}}}
\def\akt#1{\notes{\textcolor{blue}{akt: #1}}}
\def\wham#1{\smallbreak\noindent\textbf{\textit{#1}}}
\newcommand{\vu}{\vec{u}}
\newcommand{\ith}{i^\text{th}}
\newcommand{\jth}{j^\text{th}}
\def\hT{h^{(T)}}
\newcommand{\kth}{k^\text{th}}
\newcommand{\rth}{r^\text{th}}
\newcommand{\mth}{m^\text{th}}
\newcommand{\nth}{n^\text{th}}
\newcommand{\indi}{ {\small i} }
\newcommand{\indj}{ {\small j} }
\newcommand{\indm}{{\small m}}
%\newcommand{\indn}{{\small n}}
%\def\Lscr{\mathcal{L}}
\def\Lscr{\mathcal{Z}}
\def\Mscr{\mathcal{M}}
%\def\Nscr{\mathcal{N}}
%\def\Pscr{\mathcal{P}}
\def\Nscr{\mathcal{M}}
\def\Pscr{\Pi}
\def\Sscr{\mathcal{S}}
\def\R{\mathbb{R}}
\def\ie{i.e.\,}
\def\Fig#1{Fig.~\ref{#1}}
\def\Sec#1{Sec.~\ref{#1}}
\def\Subsec#1{Sec.~\ref{#1}}
\def\UZ{\underline{\mathcal{Z}}}
\def\UY{\underline{\mathcal{Y}}}
\def\uX{\underline{X}}
\def\uY{\underline{Y}}
\def\uZ{\underline{Z}}
\def\ua{\underline{\alpha}}
\def\uA{\underline{A}}
\def\upi{\underline{\pi}}
\def\ugamma{\underline{\gamma}}
\def\P{{\sf P}}
\def\Inov{I}
\def\ddx{\frac{\ud}{\ud x}}
\def\IEEEQEDclosed{\mbox{\rule[0pt]{1.3ex}{1.3ex}}}
\def\qed{\nobreak\hfill\IEEEQEDclosed}
%\def\M{\mathfrak{\bar{m}}}
%\def\N{\mathfrak{\bar{n}}}
\def\M{M}
\def\N{M}
\def\S{O}
\def\cth{c^\text{th}}
\def\lth{l^\text{th}}
\newcommand{\indicator}[1]{\mathbf{1}_{\left[ {#1} \right] }}
\newcommand{\cov}{\text{Cov}}



\newcommand{\pmat}[1]{\begin{pmatrix}#1 \end{pmatrix}}
\newcommand{\bmat}[1]{\begin{bmatrix}#1 \end{bmatrix}}


\newcommand{\vX}{X}
\newcommand{\Pos}{P}
\newcommand{\bearing}{h}
\newcommand{\bearingMean}{\hat{h}}
\newcommand{\Obs}{Z}
\newcommand{\Permm}[1]{\Pscr_{#1}}
\newcommand{\filtration}{\UZ}
\newcommand{\ProbAsso}{\beta}
\newcommand{\ProbPerm}{\gamma}
\newcommand{\assoc}[2]{#1 \equiv #2}
\newcommand{\particleX}{X}
\newcommand{\approxControlGain}{\tilde{\v}}
\newcommand{\innov}{I}

\newcommand{\velx}{\dot{x}}
\newcommand{\vely}{\dot{y}}
\def\clZ{{\cal Z}}
\def\varble{\,\cdot\,}

\def\eqdef{\mathrel{:=}}

%\newcommand{}{def}
%\newcounter{rmnum}
%\newenvironment{romannum}{\begin{list}{{\upshape (\roman{rmnum})}}{\usecounter{rmnum}
%\addtolength{\leftmargin}{-5pt}
%\setlength{\rightmargin}{4pt}
%}}{\end{list}}



% \def\RevComment#1#2{%
\newcommand{\RevComment}[3][Green]{
{\color{#1}{#2}}%
\marginpar{%
    \begin{flushleft}
	\tiny\color{#1}{#3}
    \end{flushleft}
}%
}
\newcommand{\RevChg}[2]{%
%{\color{red}{\sout{#1}}}% \xout
{\color{blue}{\uwave{#2}}}%
}
\newcommand{\RevDel}[1]{%
{\color{red}{\sout{#1}}}
}


%
\newtheorem{theorem}{Theorem}
\newtheorem{example}{Example}
\newtheorem{definition}{Definition}
\newtheorem{lemma}{Lemma}
\newtheorem{remark}{Remark}
\newtheorem{proposition}{Proposition}
\newtheorem{corollary}{Corollary}
\newtheorem{assumption}{Assumption}
% \declaretheorem[sibling=theorem]{problem}
\newtheorem{problem}{Problem}
\newcommand{\ali}[1]{\begin{align}#1\end{align}}
\def\prechi{\psi}

\def\beq{\begin{eqnarray}} 
\def\bc{\begin{center}} 
\def\be{\begin{enumerate}}
\def\bi{\begin{itemize}} 
\def\bs{\begin{small}}
\def\bS{\begin{slide}}
\def\ec{\end{center}} 
\def\ee{\end{enumerate}}
\def\ei{\end{itemize}}
\def\es{\end{small}}
\def\eS{\end{slide}}
\def\eeq{\end{eqnarray}}
\def\Xhat{\hat{X}}
\def\Yhat{\hat{Y}}
\def\Zhat{\hat{Z}}
\def\hah{\hat{h}}
\def\haphi{\hat{\phi}}
\def\hap{\hat{p}}
\def\hah{\hat{h}}
\def\haU{\widehat{U}}
\def\J{{\sf J}}
\def\Jstar{J^*}  %I didn't like a simple "J" -spm
\def\belief{\beta}
% \newcommand{\Expect}{\mathbb{E}}

\def\bfbelief{\bfmath{\beta}}
\newcommand{\newP}[1]{\noindent{\bf #1:}}
\newcommand{\Rn}{\mathbb{R}^n}
\newcommand{\px}{\frac{\partial}{\partial x}}
\newcommand{\ppt}{\frac{\partial}{\partial t}}
\newcommand{\pp}[1]{\frac{\partial}{\partial #1}}
\newcommand{\pxx}{\frac{\partial^2}{\partial x^2}}
\newcommand{\diver}{\nabla \cdot}
\newcommand{\trace}{\text{Tr}}

\newcommand{\PP}{{\sf P}}
\newcommand{\mean}{\hat{X}}
\newcommand{\ud}{\,\mathrm{d}}
\newcommand{\NGs}{N}

\def\Re{\mathbb{R}}
\def\E{{\sf E}}
\def\KL{\,\text{\rm KL}}
\def\argmin{\mathop{\text{\rm arg\,min}}}
\def\ind{\text{\rm\large 1}}
\def\varble{\,\cdot\,}

%This is p_n: \def\pxzn{p_{\text{\tiny$ X| Z^n$}}}
\def\pxz{p_{\text{\tiny$ X| Z$}}}
\def\pzx{p_{\text{\tiny$ Z | X$}}}
\def\pdzx{p_{\text{\tiny $ \ud Z | X$}}}
\def\pz{p_{\text{\tiny$Z$}}}
\def\pdz{p_{\text{\tiny$\ud Z$}}}
\def\ps{p_{\text{\tiny$S$}}}

\def\clY{{\cal Y}}

%%%

\def\Sec#1{Sec.~\ref{#1}}
\def\Thm#1{Thm.~\ref{#1}}
\def\Prop#1{Prop.~\ref{#1}}

\def\smopot{{\cal G}}

\def\bfmath#1{{\mathchoice{\mbox{\boldmath$#1$}}%
{\mbox{\boldmath$#1$}}%
{\mbox{\boldmath$\scriptstyle#1$}}%
{\mbox{\boldmath$\scriptscriptstyle#1$}}}}

\def\bfInov{\bfmath{I}}
\def\Inov{I} %So we can change symbol easily later

\def\bfPhi{\bfmath{\Phi}}
\def\bfmX{\bfmath{X}}
\def\bfmZ{\bfmath{Z}}
\def\bfmU{\bfmath{U}}
\def\bfmW{\bfmath{W}}
\def\bfxi{\bfmath{xi}}



\def\sTriangle{\hbo x{\small $\triangle$}}


\def\clB{{\cal B}}
\def\clD{{\cal D}}
\def\clE{{\cal E}}
\def\clJ{{\cal J}}
\def\clL{{\cal L}}
\def\clP{{\cal P}}
\def\clZ{{\cal Z}}
\def\clZd{{\cal Z}^{{}\, \hbox{\tiny $\blacksquare$}}}
\def\w{{\sf \Omega}}%Wong-Zakai correction term

\def\clW{{\cal W}}

\def\haphi{\hat{\phi}}
\def\hap{\hat{p}}
\def\hah{\hat{h}}
\def\haU{\widehat{U}}
\newcommand{\dthet}{\dot{q}}
\newcommand{\cm}{_{\text{\tiny CM}}}
\newcommand{\JJ}{I_{(q)}}
\newcommand{\CN}{C^{(n)}}
\newcommand{\CT}{C^{(t)}}
\newcommand{\TT}{_{t}}
\newcommand{\CC}{\mathsf{R}}
\newcommand{\fgen}{f^{\text{gen}}}
\newcommand{\bbb}{\frac{b_1 + b_2}{b+b_1 + b_2 }}
\newcommand{\C}{C_{(q)}}
\renewcommand{\Re}{\mathbb{R}}
\newcommand{\xlc}{X^{\text{LC}}}
\def\eqdef{\mathbin{:=}}
\def\Prob{{\sf P}}
\def\FRAC#1#2#3{\genfrac{}{}{}{#1}{#2}{#3}}

\def\ddx{\frac{\ud}{\ud x}}
\def\ddt{\frac{\ud}{\ud t}}
\def\J{{\sf J}}
\def\Jstar{J^*}  %I didn't like a simple "J" -spm
\newcommand{\ii}{k}

%\newcommand{\Xit}{\tilde{X}_t}
%\newcommand{\Xii}{\tilde{X}}
\newcommand{\Xii}{\bar{X}}
\newcommand{\Eii}{\tilde{E}}
%\newcommand{\hatXit}{\hat{\tilde{X}}_t}
%\newcommand{\hatXii}{\hat{\tilde{X}}}
\newcommand{\hatXii}{\bar{m}}
\newcommand{\tn}{_{t_{n}}}
\newcommand{\tnn}{_{t_{n-1}}}
\newcommand{\Sigmaii}{\bar{\Sigma}}
\def\kii{\bar{\sf K}}
\newcommand{\var}{\text{Var}}
\newcommand{\X}{X}
\newcommand{\Om}{\Omega}
\newcommand{\NN}{\mathcal{N}}

\newcommand{\val}{\tilde{\phi}}
\newcommand{\phiN}{\phi^{(N)}}
\newcommand{\kepsnorm}{{k}^{(\epsilon)}}
\newcommand{\kepsnormN}{{k}^{(\epsilon,N)}}
\newcommand{\kepsnormNN}{\bar{k}^{(\epsilon,N_2)}}
\newcommand{\keps}{{k}_{\epsilon}}
\newcommand{\kepsN}{{k}_{\epsilon}^{(N)}}
\newcommand{\Keps}{{K}_\epsilon}
\newcommand{\KepsN}{{K}_\epsilon^{(N)}}
\newcommand{\Kepsnorm}{\bar{K}_\epsilon}
\newcommand{\KepsnormN}{\bar{K}_\epsilon^{(N)}}
\newcommand{\phiepsN}{{\phi}^{(N)}_\epsilon}
\newcommand{\phieps}{{\phi}_{\epsilon}}
\newcommand{\LepsN}{{L}_\epsilon^{(N)}}
\newcommand{\HH}{H^1_0}
\newcommand{\sumiN}{\sum_{i=1}^N}
\newcommand{\phiM}{\phi^{(M)}}
\newcommand{\phiMN}{\phi^{(M,N)}}
\newcommand{\Deltaeps}{{\Delta_\rho}^{(\epsilon)}}
\newcommand{\DeltaepsN}{{\Delta_\rho}^{(\epsilon,N)}}
\newcommand{\rhoepsN}{\rho^{(\epsilon,N)}}
\newcommand{\rhoeps}{\rho^{(\epsilon)}}
\newcommand{\Teps}{{T}_{\epsilon}}
\newcommand{\Aeps}{{A}^{(\epsilon)}}
\newcommand{\TepsN}{{{T}^{(N)}_\epsilon}}
\newcommand{\mueps}{{\mu}^{(\epsilon)}}
\newcommand{\geps}{{g}_{\epsilon}}
\newcommand{\pr}{\rho}
\newcommand{\neps}{{n}^{(\epsilon)}}
\newcommand{\Ltwoeps}{{L^2_0(\mueps)}}
\newcommand{\FF}{C_0(\Omega)}
\newcommand{\infnorm}[1]{\|#1\|_\infty}
\newcommand{\PhiepsN}{{\Phi}^{(\epsilon,N)}}
\newcommand{\vepsN}{{{\sf K}}_\epsilon^{(N)}}
%\floatname{algorithm}{Procedure}
%\renewcommand{\algorithmicrequire}{\textbf{Input:}}
%\renewcommand{\algorithmicensure}{\textbf{Output:}}
\newcommand{\Geps}{{G}_\epsilon}
\newcommand{\GepsN}{{G}_\epsilon^{(N)}}
\newcommand{\Tm}{{\bf T}}
\newcommand{\Km}{{\bf K}}
\newcommand{\hm}{{\bf h}}
\newcommand{\epsbar}{\bar{\epsilon}}
\newcommand{\meas}{\mu}
\newcommand{\aht}[1]{{\color{red} \text{#1}}}
\newcommand{\mN}{m^{(N)}}
\newcommand{\SigN}{\Sigma^{(N)}}
\newcommand{\GammaN}{\Gamma^{(N)}}
\newcommand{\OmegaN}{\Omega^{(N)}}
\newcommand{\snorm}[1]{\|#1\|_s}
\newcommand{\Fnorm}[1]{\|#1\|_F}
\newcommand{\Xs}{\bar{X}}
\newcounter{rmnum}
\newenvironment{romannum}{\begin{list}{{\upshape (\roman{rmnum})}}{\usecounter{rmnum}
\setlength{\leftmargin}{14pt}
\setlength{\rightmargin}{8pt}
\setlength{\itemsep}{2pt}
\setlength{\itemindent}{-1pt}
}}{\end{list}}

\newcounter{anum}
\newenvironment{alphanum}{\begin{list}{{\upshape (\alph{anum})}}{\usecounter{anum}
\setlength{\leftmargin}{14pt}
\setlength{\rightmargin}{8pt}
\setlength{\itemsep}{2pt}
\setlength{\itemindent}{-1pt}
}}{\end{list}}
\newcommand{\Ten}{{\sf T}}
\newcommand{\preps}{\pr_\epsilon}
\newcommand{\barphieps}{\bar{\phi}_\epsilon}
\newcommand{\barphiepsN}{\bar{\phi}_\epsilon^{(N)}}
\newcommand{\barnepsN}{\bar{n}_\epsilon^{(N)}}
\newcommand{\nepsN}{{n}_\epsilon^{(N)}}
\newcommand{\barkeps}{\bar{k}_\epsilon}
\newcommand{\barphivec}{\bar{{\sf  \Phi }}}
\newcommand{\barTen}{\bar{\sf T}}
\newcommand{\hvec}{{\sf h}}
\newcommand{\phivec}{{\sf \Phi}}

\def\sTriangle{\hbox{\small $\triangle$}}
%%%% Chi's new definitions

\def\phiM{\phi^{M}}
\def\Ten{{\sf T}}
\newcommand{\GN}{G^{(N)} }
\newcommand{\sRicc}{\sqrt{\text{Ricc}}}
\newcommand{\Ttn}{T_{\frac{\epsilon}{n}}^n}
\newcommand{\Pt}{e^{\epsilon\Delta_\pr}}
\newcommand{\calX}{\mathcal{X}}
\newcommand{\calB}{\mathcal{B}}
\newcommand{\calS}{\mathcal{S}}
\newcommand{\calA}{\mathcal{A}}
%\newcommand{\calT}{\mathcal{T}}
%\newcommand{N}{{T}^{(\epsilon,N_2)}}
%\newcommand{\Deltaeps}{\frac{1}{\epsilon}(I -e^{-\epsilon\Delta_\rho})}
\newcommand{\Alg}{optimal transport particle filter}
\newcommand{\ALG}{{OTPF}}

\newcommand{\norm}[1]{\left\|#1\right\|}

\newcommand{\Lip}{\text{Lip}}

\title{\LARGE \bf A Time-Reversal Control Synthesis for Steering the State of \\Stochastic Systems
}

%\author{ \parbox{3 in}{\centering Huibert Kwakernaak*
		%         \thanks{*Use the $\backslash$thanks command to put information here}\\
		%         Faculty of Electrical Engineering, Mathematics and Computer Science\\
		%         University of Twente\\
		%         7500 AE Enschede, The Netherlands\\
		%         {\tt\small h.kwakernaak@autsubmit.com}}
	%         \hspace*{ 0.5 in}
	%         \parbox{3 in}{ \centering Pradeep Misra**
		%         \thanks{**The footnote marks may be inserted manually}\\
		%        Department of Electrical Engineering \\
		%         Wright State University\\
		%         Dayton, OH 45435, USA\\
		%         {\tt\small pmisra@cs.wright.edu}}
	%}


 \author{Yuhang Mei$^{\star}$, Amirhossein Taghvaei$^{\star}$, Ali Pakniyat$^{\dagger}$
 \thanks{Yuhang Mei and Amirhossein Taghvaei are supported by the National Science Foundation (NSF) award EPCN-2318977 and EPCN-2347358.}
     \thanks{$^\star$Department of Aeronautics \& Astronautics, University of Washington, Seattle; {\tt\small yuhangm@uw.edu,amirtag@uw.edu}.}   
	\thanks{$^\dagger$Department of Mechanical Engineering,
University of Alabama; {\tt\small apakniyat@ua.edu}.}
}



\newcommand{\Ali}[1]{{\color{magenta}[Ali: #1]}}

\begin{document}
	
	
	
      \maketitle
	 \thispagestyle{empty}
	 \pagestyle{empty}
	
	
	%%%%%%%%%%%%%%%%%%%%%%%%%%%%%%%%%%%%%%%%%%%%%%%%%%%%%%%%%%%%%%%%%%%%%%%%%%%%%%%%
    
        



        
	 \begin{abstract}
This paper presents a novel approach for steering the state of a stochastic control-affine nonlinear system to a desired target within a finite time horizon. Our method leverages the time-reversal of diffusion processes to construct the required feedback control law. Specifically, the control law is the so-called score function associated with the time-reversal of random state trajectories that are initialized at the target state and are simulated backwards in time.  A neural network is trained to approximate the score function, enabling applicability to both linear and nonlinear stochastic systems.  Numerical experiments demonstrate the effectiveness of the proposed method across several benchmark examples. 
%Furthermore, we provide an error analysis for linear systems with Gaussian noise while extending the analysis to more general settings remains a topic for future research.


	\end{abstract}
	%

   
        

 
    \section{Introduction}
Steering the state of a stochastic system to a target state or distribution is a fundamental problem in stochastic control and useful in applications such as stochastic thermodynamics~\cite{peliti2021stochastic,chen2019stochastic,fu2021maximal}, machine learning~\cite{chen2023generative,domingo2024stochastic,rapakoulias2024go}, and robotics~\cite{shah2012stochastic,gorodetsky2018high,theodorou2011iterative,okamoto2019optimal,liu2024optimal}. 
One prominent framework for studying this problem is the  theory of Schr\"odinger bridges for diffusion processes~\cite{wakolbinger1990schrodinger,pavon1991free}, which aims to find an optimal control law that drives the system from a given initial distribution to a specified target distribution over a finite horizon. Exact solutions are available for linear stochastic systems with Gaussian initial and target distributions~\cite{chen2015optimal,chen2015part2,teter2024schr,teter2025markov}, and numerical procedures extend these results to non-Gaussian~\cite{liu2023generalized,mei2024flow} and nonlinear drift~\cite{caluya2021wasserstein} cases. 
Beyond Schr\"odinger bridge theory, alternative approaches using the stochastic maximum principle and convex duality have been proposed to derive optimal control policies that steer the state toward desired target  points~\cite{pakniyat2021steering, pakniyat2021partially} or distributions~\cite{pakniyat2022convex, pakniyat2024distributionally}.
{However, the reliance on optimality in both these and Schr\"odinger-based approaches makes controller identification for nonlinear settings computationally expensive, as it requires solving complex partial differential equations.}
% \sout{However, both these and Schr\"odinger-based approaches are restricted to linear stochastic dynamic models, and extension to nonlinear settings require solution to computationally expensive partial differential equations}. 
% \Ali{Note: my results hold for general nonlinear systems, but their solution for nonlinear systems are challenging. Also \cite{caluya2021wasserstein} is a nonlinear Schr\"odinger bridge result.)}

In contrast, this paper shifts attention away from optimality (and initial-state dependence) toward designing a feedback law that guarantees finite-time attractability of a chosen target. Specifically, we consider a stochastic process $X$ that is governed by a control-affine nonlinear stochastic system~\eqref{eq:sde} and propose a framework for deriving feedback control laws that steer the state $X_t$ toward a target state $X_T=x_f$, at a finite time horizon $T$, without requiring optimality relative to the initial condition.

\begin{figure}[t]
\centering
\includegraphics[width=0.98\hsize]{figs/schematic-drawing.pdf}
\caption{Illustration of the proposed time-reversal methodology to steer the state $X_t$ from $x_0$ to $x_f$. The auxiliary process $Z_t$ is simulated backward in time from $x_f$. The dynamics of the time-reversal $\tilde Z_t:=Z_{T-t}$  is forward in time, with an absorption property to $x_f$. This dynamics is used to design the control law $U_t=k(t,X_t)$. The expressions for  the functions $h$, $\tilde h$, and $k$ appear in Section~\ref{sec:methodology}.}
\label{fig:schematic}
\end{figure}

Our approach is inspired by time-reversal theory of diffusions~\cite{anderson1982reverse,haussmann1986time,cattiaux2023time}, which has recently gained attention in machine learning through diffusion-based generative models for images~\cite{NEURIPS2020_4c5bcfec,song2020denoising,song2019generative,songscore2021score}. In these models, one first simulates a stochastic differential equation (SDE) that gradually adds noise to transform complex data (e.g., images) into a simpler, typically Gaussian distribution. The trajectory of this noising process is then used to learn the so-called score function, which in turn is used to run the SDE in reverse. Sampling from the Gaussian and applying this reverse process (the ``denoising'' procedure) generates new samples resembling the original data. From a control-theoretic viewpoint, the learned score function serves as a control law that transforms the Gaussian distribution to the target data distribution. A key benefit of this procedure is its computational tractability as it involves solving a regression problem rather than resorting to dynamic programming or the maximum principle, though at the cost of losing optimality.

Inspired by these diffusion generative models, we draw an analogy for our setting by simulating the state of stochastic system backward in time from a given target state $x_f$ (or a small Gaussian distribution around it). Reversing this backward process in the forward-time direction ensures convergence to the target, with the learned score function serving as the feedback control law. To provide an intuitive understanding of our approach, we outline the key steps involved in constructing the proposed feedback control law, accompanied with an illustration in Figure~\ref{fig:schematic}. 
% the complex data distribution is replaced by the target state, and the score function serves as the control law to achieve the task.  By constructing a feedback control law associated with the score function, we aim to steer the system toward the target state. In other words, by modifying the drift term of SDE to enforce the desired terminal conditions based on the time-reversal of a stochastic process that starts from the target state, our objective is to obtain an SDE for the control system, and hence obtain a control law which steers the trajectories to the desired terminal state as follows. 

% A special, yet crucial, case of Schr\^odinger bridge problem is when the final distribution is a Dirac delta distribution, i.e. the goal of steering  the state of a stochastic system to a given target state, in finite time, motivated by applications such as landing reusable rockets. While this problem is seemingly a simpler problem compared to 



% Stochastic control is central to many applications, including finance \cite{pham2009continuous, fleming2004stochastic}, generative models \cite{chen2023generative,rombach2022high,ramesh2022hierarchical,domingo2024stochastic}, stochastic thermodynamics \cite{Sekimoto2010,peliti2021stochastic}, and robotics .

% A key challenge within stochastic control is steering the state of a stochastic system to a target state or distribution, a task crucial in applications such as landing reusable rockets, rarefied gas dynamics, and interacting gas models. 



% Several methods have been developed to address this challenge. Namely, Schr$\ddot o$dinger bridge theory has been employed to steer a linear stochastic systems to a desired Gaussian distribution over a finite time horizon, \cite{chen2015optimal,chen2018steering}.  An extension to non-Gaussian distributions has also been explored for linear systems \cite{mei2024flow} at the expense of losing optimality. However, the applicability of these methods remains limited to linear systems. 

% For nonlinear systems, optimal control strategies have been explored. The maximum principle has been used to derive optimal control laws that steer the state towards a desired target state \cite{pakniyat2021steering} but this methodology does not accommodate steering towards desired probability distributions. The steering of nonlinear systems to target distributions has been examined in \cite{pakniyat2022convex} by formulating the optimal control problem in the measure space, though this approach requires solving computationally expensive optimizations over families of partial differential equations (PDEs).


% In contrast to these works, in this paper, we propose a time-reversal method that provides a framework for obtaining a feedback control law capable of steering any control affine nonlinear stochastic system to the target, though without guaranteeing optimality.

% Time-reversal has recently gained attention as a method for denoising diffusion generative models \cite{NEURIPS2020_4c5bcfec,song2020denoising,song2019generative,songscore2021score}. In these works, a stochastic differential equation (SDE) is constructed to transform a complex data distribution (e.g., an image) to a known prior distribution, typically Gaussian. Then, the time-reversal of this SDE is applied to remove noise and generate images. Inspired by this approach, we draw an analogy where the target state corresponds to the complex data distribution. By constructing a feedback control law associated with the score function, we aim to steer the system toward the target state. In other words, by modifying the drift term of SDE to enforce the desired terminal conditions based on the time-reversal of a stochastic process that starts from the target state, our objective is to obtain an SDE for the control system, and hence obtain a control law which steers the trajectories to the desired terminal state as follows. 

% To provide an intuitive understanding of our approach, we outline the key steps involved in constructing a feedback control law using time-reversal principles, as follows:
\begin{enumerate}
    \item \textbf{Auxiliary Process $\boldsymbol{Z}$:} We construct an auxiliary process $Z$, initialized at the desired terminal state $x_f$. This process generates a probability density function $p(t,x)$ with a delta distribution at the desired terminal location at its initial time, i.e., $p(0,x) = \delta_{x_f}(x)$ while producing spread-out distributions at other times. This serves as the foundation for defining a time-reversed process.

    \item \textbf{Time-Reversed Auxiliary Process $\boldsymbol{\tilde{Z}}$}: By reversing time for the auxiliary process $Z$, we obtain a new process $\tilde{Z}_t =Z_{T-t}$ with the probability density function $\tilde{p}(t,x) = p(T-t,x)$ possessing a delta distribution at the desired terminal location (matching the target state) at its terminal time, i.e., $\tilde{p}(T,x) = \delta_{x_f}(x)$. The governing dynamics of $\tilde{Z}$ possess an absorption property for the desired terminal state, i.e., almost all sample paths are attracted to and absorbed by the desired state. This property is critical for ensuring that trajectories converge to the target state.

    \item \textbf{Construction of Feedback Control Law for $\boldsymbol{X}$:} The feedback law synthesis for $X$ leverages the absorption property of the time-reversed process $\tilde{Z}$, to yield closed loop dynamics that enforce almost sure convergence to the target state within a fixed time horizon. Crucially, if the initial condition $x_0$ of the original process $X$ lies within the support of the initial distribution of $\tilde{Z}$, then $x_0$ belongs to the fixed-time stochastic region of attraction of the desired terminal state almost surely.
\end{enumerate}

\noindent
{\bf Relation to diffusion bridges:} A diffusion bridge is defined as the law of the diffusion processes $X$ conditioned on the event that $X_0=x_0$ and $X_T=x_f$. Using Doob's $h$-transform \cite[Sec. 7.5]{sarkka2019applied}, the conditioned process satisfies the same SDE as the original diffusion but with an added drift that can be interpreted as a control law steering the state to the target. However, this control law is distinct from the one we propose in this paper. In particular, our control law is not generally of gradient form, whereas the Doob $h$-transform yields the gradient of the log-transition density.  Interestingly, the two control laws coincide in the linear Gaussian setting, suggesting a deeper connection between the approaches. From a computational standpoint, there is a recent work~\cite{10.1093/biomet/asaf048} that applies time-reversal and score matching to simulate diffusion bridges. However, their proposed method  requires applying time reversal and score-function approximation twice~\cite[Sec. 2.3]{10.1093/biomet/asaf048}, with potential error accumulation, while our method requires only a single such operation.
A detailed numerical comparison to fully characterize the trade-offs is planned for future investigation.

 
 % Our method is closely related to the recent work~\cite{10.1093/biomet/asaf048}, which applies time-reversal and score matching to simulate the law of the diffusion processes $X$ conditioned on the event that $X_T=x_f$. Using Doob's $h$-transform \cite[Sec. 7.5]{sarkka2019applied} the conditioned process satisfies the same SDE as the original diffusion but with an added drift that acts as a feedback control with an elegant optimality interpretation~\cite{chen2016relation}. Both this control law and the one derived from our method guarantee convergence to the target, though their forms differ. In particular, our control law is not generally of gradient form, whereas the Doob $h$-transform yields the gradient of the log-transition density. From a computational standpoint, approximating the Doob $h$-transform is more demanding: it requires applying time reversal and score-function approximation twice~\cite[Sec. 2.3]{10.1093/biomet/asaf048}, with potential error accumulation, while our method requires only a single such operation. Interestingly, the two control laws coincide in the linear Gaussian setting, suggesting a deeper connection between the approaches. A detailed numerical comparison will be pursued in future work. 

% A stochastic process conditioned to start and end at two given states within a fixed time interval is known as a bridge. The way to realize such unique bridge is through Doob's $h$-transform \cite[p. 83]{rogers2000diffusions} which modifies the drift term using the gradient of the logarithm of the transition density to the terminal state. A closely related work for simulating the diffusion bridge was recently proposed in \cite{10.1093/biomet/asaf048}. 
% %The key distinction is that their method simulates the time-reversed conditioned dynamics, whereas our method starts with the auxiliary process $Z$ and constructs a feedback control law for the forward process $X$. 
% From the control perspective, Doob's $h$-transform term can itself be interpreted as a feedback control law. However, as discussed in \cite[Sec. 2.3]{10.1093/biomet/asaf048}, obtaining this feedback control law to simulate the forward conditioned process requires approximating both the score of the unconditioned forward process and the score of the diffusion bridge, which introduces error accumulation. Interestingly, in the general setting, our feedback control law is not equivalent to Doob's $h$-transform formulation in \cite[Eq. (2)]{10.1093/biomet/asaf048}; yet in the linear Gaussian case, the two approaches coincide exactly, and the analysis of the linear Gaussian setting appears in Section \ref{sec:lin-gau}. This observation suggests a potential conceptual connection between the two approaches. A detailed numerical comparison of the two methods will be carried out as part of future work. 
% In this paper, although we construct a bridge between two fixed states through the stochastic system, the resulting feedback control law is not equivalent to the classical Doob's $h$-transform \cite[p. 83]{rogers2000diffusions}. A closely related line of work was recently proposed in \cite{10.1093/biomet/asaf048}, where diffusion bridges are simulated by approximating the time-reversed dynamics. Compared to \cite{10.1093/biomet/asaf048}, the distinct is that our approach leverages the score function of auxiliary process $Z$ and synthesis feedback control law through time-reversal of process $Z$ conditioned on two fixed terminal states. Interestingly, in the linear Gaussian setting we find that our feedback controller coincides exactly with the Doob's $h$-transform conditioned on two fixed terminal states. This highlights both the conceptual connection between the two approaches and the novel control-theoretic perspective introduced in our work.
% In another recent paper \cite{10.1093/biomet/asaf048}, a related approach was proposed for simulating diffusion bridges, where the time-reversal representation of diffusions is combined with score-matching techniques to learn the score function and efficiently generate conditioned sample paths.

The paper is organized as follows. Section \ref{sec:problem_formulation} presents the problem formulation and the necessary background on the time-reversal theory of diffusions. Section \ref{sec:methodology} presents our proposed  time-reversal methodology for control synthesis, along with the analysis of the linear Gaussian setting. Section \ref{sec:results} presents numerical experiments on several benchmark examples, demonstrating the performance of the proposed method in both linear and nonlinear settings. 

% \at{Suggested outline: importance of stochastic control. Importance of steering a stochastic system to a desired target. The existing methods and their limitations. Your proposed approach. Talk about time-reversal. And how you use it to solve the problem. Paper's outline. }

\subsection{Notation}
% \at{introduce a notation section here. Intoduce the space of functions $H^1(\rho)$. Look at the notation section here: https://arxiv.org/pdf/1902.07263}
% The space of measurable functions $f: \R^d \rightarrow \R$ such that $\|f\|_{L^p(\rho)} \mathrel{:=} (\int|f(x)|^p \rho(x) \ud x)^{1/p} < \infty$ is denoted as $L^p(\rho)$. The space $H^1(\rho)$ is the space of functions $f\in L^2(\rho)$ whose derivative (defined in the weak sense) is in $L^2(\rho)$. 
The notation $\frac{\partial}{\partial x}$ is used to denote the derivative with respect to the $x$ variable. For example, for a smooth function $f:\R^n \to \Re$,  $\frac{\partial f}{\partial x}:\R^n\to \R^n$ denotes the gradient of $f$ with respect to $x$. And for a smooth vector-field $f:\R^n \to \Re^n$,  $\frac{\partial f}{\partial x}:\R^n\to \R^{n\times n}$ denotes the Jacobian.  The probability density function (PDF) of a multivariate Gaussian distribution with mean vector  $m$ and covariance matrix $\Sigma$ is denoted by $\mathcal{N}(x;\mu,\Sigma)$. For a positive definite matrix $P$, the weighted $2$-norm of a vector $x \in \R^n$, is denoted by $\|x\|_P$, and defined as $\|x\|_P := \sqrt{x^\top Px}$.





    
    \section{Problem Formulation and Background}\label{sec:problem_formulation}
\subsection{Problem setup}
Consider a  control system governed by a control-affine stochastic differential equation (SDE)
\begin{equation}\label{eq:sde}
    \ud X_t = f(X_t)\ud t + g(X_t) (U_t \ud t + \epsilon \ud W_t), ~~X_0 = x_0,
\end{equation}
where $X_t \in \mathbb{R}^n$ is the state, $U_t \in \R^m$ is the control input, $W_t \in \R^m$ is the standard Wiener process that denotes the process noise, $f:\R^n\to\R^n$ and $g:\Re^n \to \Re^{n\times m}$ are drift and diffusion functions, and $\epsilon >0$ is a positive parameter that denotes the strength of the noise. Let $\mathcal{F}_t = \sigma(W_s; 0\leq s \leq t)$ be the filtration generated by the Wiener process. The control input $U_t$ is constrained to be adapted to the filtration $\mathcal{F}_t$. We are interested in the problem of designing the control input that steers the state of the system to a given target state. 

 % \begin{remark}
 %     In this paper, we only consider control affine systems, and assume the noise enters the system through the same channel as control.
 % \end{remark}
% We are interested in solving the following stochastic control problem.

\begin{problem}\label{prob:steer-to-y}
    Given a target state $x_f \in  \R^n$ and a fixed terminal time $T>0$, find a feedback control law $\{U_t=k(t,X_t);t\in[0,T]\}$, for a function $k: [0,T] \times \R^n \rightarrow \R^m$, such that  $X_T = x_f$ almost surely.
\end{problem}

Note that the $\mathcal{F}_t$-adaptability condition is automatically satisfied when $U_t$ is expressed as a function of $X_t$.  We also consider a relaxed version of the problem, where the  equality is relaxed to a bound on the average distance from the target. This relaxation allows us to obtain control laws that are not singular as $t\to T$. See Section~\ref{sec:singular}.

\begin{problem}\label{prob:approx-steer-to-y}
    Given a target state $x_f \in  \R^n$, a fixed terminal time $T>0$, and error tolerance $\delta>0$, find a feedback control law $\{U_t=k(t,X_t);t\in[0,T]\}$, for a function $k: [0,T] \times \R^n \rightarrow \R^m$, such that  $\mathbb E[\|X_T - x_f\|^2]\leq \delta$.
\end{problem} 

\begin{remark}[Modelling assumptions]
   The control-affine structure of the dynamic model in~\eqref{eq:sde} is critical for the applicability of the proposed method. This model is common in robotics, aerospace, and finance \cite{zhou2023safe,liu2019fuel,calafiore2008multi}. The time-invariant assumption for the functions $f$ and $g$ is made for ease of presentation and can be relaxed to be time-varying. 
   %And the number $\epsilon$  can be replaced by a matrix. 
\end{remark}

% \subsection{Outline of the Proposed Method}

Our solution methodology for these problems is based on the time-reversal of diffusion theory, which we review  next. 
% In other words, the problem is to construct a stochastic bridge. 


\subsection{Time-Reversal of Diffusions}
\label{sec:tr}
Consider a diffusion process $Z \mathrel{:=} \{Z_t\in \R^n;0 \leq t \leq T\}$ governed by the following SDE 
\begin{equation}\label{eq:fsde}
    \ud Z_t \!=\! h(Z_t) \ud t + \epsilon g(Z_t) \ud W_t,\quad Z_0=z,
    \end{equation}
    where   $h:  \R^n \rightarrow \R^n$ is a drift function, and $z$ is an initial condition, that are later designed as part of our methodology, in Section~\ref{sec:naive-alg}.  We make the following assumption about the functions $h$ and $g$. 
    
\begin{assumption}\label{asum-fg}
    The functions $h$ and $g$ are smooth and globally Lipschitz. 
\end{assumption}

Assumption \ref{asum-fg} implies that SDE \eqref{eq:fsde} admits a unique strong solution \cite[Thm. 5.2.1.]{oksendal2003stochastic}. This allows us to define the time-reversal of $Z_t$ according to 
\begin{equation*}
    \tilde{Z} \mathrel{:=} \{\tilde{Z}_t =Z_{T-t} ;\, 0 \leq t \leq T\}. 
\end{equation*} 
The time-reversal theory is concerned with obtaining the SDE for the reversed process $\tilde Z_t$. In order to do so, we follow the approach presented in~\cite{haussmann1986time} which, in addition to Assumption~\ref{asum-fg}, requires a suitable integrability condition on the density of $Z_t$. In particular, letting $p(t,x)$ denote the probability density function of $Z_t$, it is required that 
    \begin{equation}\label{eq:asum-bound}
        \int_{t_0}^T \int_\Omega |p(t,x)|^2 + \|g(x)^\top \frac{\partial p}{\partial x} (t,x)\|^2 \ud x \ud t < \infty ,
    \end{equation}
    for any open bounded set $\Omega \subset \Re^n$  and $t_0>0$.
%     where $\partial_x p_t(x)$ denotes the partial derivative of $p_t(x)$ with respect to $x$. 
This condition is valid under  the following assumption about the drift and diffusion functions. 

% This SDE is usually stated under the assumption that the diffusion function is non-singular, i.e. $G(t,x):=g(t,x)g(t,x)^\top$ is  uniformly positive-definite for all $(t,x)\in [0,T]\times \R^n$~\cite{anderson1982reverse,cattiaux2023time}. However, this assumption is restrictive due to the role that the function $g$ plays in the stochastic control problem~\eqref{eq:sde}. In particular, the columns of $g$ are the directions that the control input affects the state. Therefore, a positive-definite assumption on $G$ implies full control authority which is restrictive in many control applications. Instead, we follow the alternative approach for constructing the SDE of the reversed process in \cite{haussmann1986time} which requires the following weaker assumption about $g$. 

\begin{assumption}\label{asum-brackets}
    For all $x\in \R^n$, the subspace generated by the Lie-brackets\footnote{The Lie-bracket of two smooth vector-fields $g_i:\R^n \to \R^n$ and $g_j:\R^n\to \R^n$ is defined as $[g_i,g_j](x)=\frac{\partial g_j}{\partial x}(x)g_i(x) - \frac{\partial g_i}{\partial x}(x)g_j(x)$} of the elements of $\{h(x),g_1(x),\ldots,g_m(x)\}$, with $g_i$ the $i$-th column of~$g$, spans the entire space $\R^n$. 
\end{assumption}

Assumption \ref{asum-brackets} is knwon as the H\"ormander condition that ensures the generator associated with the diffusion process~\eqref{eq:fsde} is hypoelliptic, implying that $Z_t$ admits a smooth density $p(t,x)$ for any $t>0$~\cite{chaleyat1984hypoellipticity,hormander1967hypoelliptic}. As a result, condition~\eqref{eq:asum-bound} is satisfied. With a smooth density, one can define the so-called score function $s:(0,T] \times \Re^n \to \Re^n$ whose $i$-th component is given by
  \begin{align}\label{eq:score-def}
            s_i(t,x) \mathrel{:=} \frac{1}{p(t,x)} \sum_{j=1}^n  \frac{\partial}{\partial x_j}(G_{i,j}(x)p(t,x)),
            %\quad %\forall (t,x)\in (0,T]\times \R^n 
        \end{align}
        where $G_{i,j}(x)$ is the $(i,j)$-entry of the matrix $G(x):=g(x)g(x)^\top \in \R^{n\times n}$. 
Then, according to \cite[Thm. 2.1]{haussmann1986time},
the reversed process $\tilde Z_t$ satisfies the SDE
\begin{align}
            \ud \tilde{Z}_t = -h(\tilde Z_t)\ud t +  \epsilon^2 s(T-t,\tilde Z_t) \ud t+ \epsilon g(\tilde{Z}_t) \ud  \tilde W_t,    \label{eq:tr-sde} 
        \end{align}
        where $\tilde W_t \in \R^m$ is the standard $m$-dimensional Wiener process. Note that the reversed process satisfies the condition $\tilde Z_T=Z_0 = z$. Therefore, the dynamics generated by the score function and  the drift term has an absorption property for the state $z$ at terminal time $T$. This is the basis for our methodology in Section~\ref{sec:naive-alg}. 
        % and the new function $\tilde h$ is related to $h$ and score function $s$ with the identity
        % \begin{align}\label{eq:tilde-h}
        %     \tilde h(T-t,x) + h(t,x) = \epsilon^2 s(t,x).
        % \end{align}
      
         \begin{remark}
        The SDE for the time-reversal process may also be obtained using Girsanov theorem, as in~\cite{cattiaux2023time}, but this alternative approach requires the diffusion function to be non-singular, i.e. $G(x):=g(x)g(x)^\top$ be uniformly positive-definite for all $x\in \R^n$~\cite{anderson1982reverse,cattiaux2023time}. This assumption is restrictive due to the role that the function $g$ plays in the stochastic control problem~\eqref{eq:sde}. In particular, the columns of $g$ are the directions that the control input affects the state. Therefore, a positive-definite assumption on $G$ implies full control authority which is restrictive in many control applications. Assumption~\ref{asum-brackets} is a weaker assumption that allows the construction of the time-reversal SDE for a degenerate diffusion. It is also in agreement with the local accessibility condition in geometric control theory  for the deterministic version of the system~\eqref{eq:sde},  when $\epsilon=0$~\cite{krener1974generalization,sussmann1987general}. 
    \end{remark}
    
    The formula for the score function~\eqref{eq:score-def} depends on the density $p(t,x)$ which is explicitly available only in the special linear Gaussian setting~(see Section~\ref{sec:lin-gau}). In a general nonlinear and non-Gaussian setting, the score function  is approximated as the solution to a stochastic optimization problem, as described in the next subsection. 
% \Ali{It would be beneficial to add a sentence about the score function itself, before moving to the next section regarding alternative approaches for its numerical determination.}  
\subsection{Score function approximation}
It is numerically useful to note that the score function is the solution to the minimization problem $\min_\psi\,J(\psi)$ where 
 \begin{align}\label{eq:score-opt}
       J(\psi):= \mathbb{E}[ \frac{1}{2}\|\psi(t,Z_t)\|^2 + \sum_{i,j=1}^n G_{i,j}(Z_t) \frac{\partial \psi_i}{\partial x_j} (t,Z_t)) ],
    \end{align}
    and the expectation is both over $t \sim \text{Unif}[0,T]$ and $Z_t$, solution to~\eqref{eq:fsde}. 
    This follows by writing the expectation as the integral with respect to the density $p(t,x)$, application of integration by parts on the second term,  and expressing the objective function~\eqref{eq:score-opt} as 
    \begin{align*}
        \mathbb{E}\left[ \frac{1}{2}\|\psi(t,Z_t)-s(t,Z_t)\|^2   \right] + \text{(constant)},
    \end{align*}
    where the constant is independent of $\psi$~\cite[Thm. 1]{hyvarinen2005estimation}. The optimization \eqref{eq:score-opt} is known as implicit score matching \cite{tang2024score}. This optimization procedure is later used in the construction of our numerical algorithm in Section~\ref{sec:Algorithm}.  

% Numerically, the score function is approximated by sampling $\t^i\sim \text{Unif}[0,T]$ and $Z^i_t$ by simulating~\eqref{eq:fsde}, approximating the expectation in the objective function~\eqref{eq:score-opt} in terms of the samples $(t_i,Z^i_t)_{i=1}^N$, and optimizing $\psi$ over a parametric class of functions, e.g. neural networks. 
   

    \subsection{Probability transition kernels}\label{sec:kernel}
   This section introduces notations and definitions for probability transition kernels associated with SDEs~\eqref{eq:fsde} and~\eqref{eq:tr-sde} that are later used in the proof of Propostions~\ref{prop:sol2p1} and~\ref{prop:terminal-error}. 
    Let $\kappa_{t,s}(x'|x)$ denote the probability transition kernel from time $s$ to time $t$ that is associated with SDE~\eqref{eq:fsde},
i.e. the conditional probability density function of $Z_t=x'$ given $Z_s=x$, for any $t\geq s>0$. Similarly, let $\tilde \kappa_{t,s}(x'|x)$ denote the probability transition kernel for SDE~\eqref{eq:tr-sde}. Using the probability transition kernel, the joint probability density function of $(Z_t,Z_s)$ satisfies
% \Ali{It could be beneficial to start the argument by $ \int\limits_{B_{x} \times B_{x'}} P_{Z_t,Z_s}(x',x) \ud x \ud x' = \int\limits_{B_{x} \times B_{x'}} \kappa_{t,s}(x'|x) p(s,x) \ud x \ud x'$ for arbitrary Borel sets $B_{x}, B_{x'} \subset \R^n$, and then express the following (just to avoid potential confusions of working with zero probabilities).}
\begin{align*}
\int\limits_{B_{x} \times B_{x'}} \hspace{-9pt}P_{Z_t,Z_s}(x',x) \ud x^\top \! \ud x'\! =\! \int\limits_{B_{x} \times B_{x'}} \hspace{-9pt}\kappa_{t,s}(x'|x) p(s,x) \ud x^\top \!\ud x'
    % P_{Z_t,Z_s}(x',x) =  \kappa_{t,s}(x'|x) p(s,x).
\end{align*}
for arbitrary Borel sets $B_{x}, B_{x'} \subset \R^n$, implying   
\begin{align*}
     P_{Z_t,Z_s}(x',x) =  \kappa_{t,s}(x'|x) p(s,x).
\end{align*}
Similarly, for $(\tilde Z_t,\tilde Z_s)$ we have 
\begin{align*}
    P_{\tilde Z_t,\tilde Z_s}(x',x) =  \tilde \kappa_{t,s}(x'|x) p(T-s,x), 
\end{align*}
% \begin{align*}
% \int\limits_{B_{x} \times B_{x'}} P_{\tilde Z_t,\tilde Z_s}(x',x) \ud x \ud x' = \int\limits_{B_{x} \times B_{x'}} \tilde \kappa_{t,s}(x'|x) p(T-s,x) \ud x \ud x'
%     % P_{Z_t,Z_s}(x',x) =  \kappa_{t,s}(x'|x) p(s,x).
% \end{align*}
where we used the fact that $\tilde Z_s = Z_{T-s}$ with probability density $p(T-s,x)$. The identity $(Z_T,Z_{T-t})= (\tilde Z_0,\tilde Z_t)$ implies the equality $P_{Z_T,Z_{T-t}}(x',x)= P_{\tilde Z_{t},\tilde Z_0}(x,x')$, concluding the relationship 
\begin{align}\label{eq:kappa-identity}
   &\kappa_{T,T-t}(x'|x) p(T-t,x) = \tilde \kappa_{t,0}(x|x') p(T,x')
\end{align}
between the two kernels $\kappa$ and $\tilde \kappa$.


 \section{Proposed Methodology}
 \label{sec:methodology}



    % $b(t,x)$ is known as F$\Ddot{\text{o}}$llmer's drift term or score function \cite{follmer2005entropy,hyvarinen2005estimation}.
% where the minimization is over all func
% Let $p_t(x)$ denote the probability density function of $X_t$. We define the time-reversal stochastic process 
% \begin{equation}\label{eq:tr-def}
%     \tilde{X} \mathrel{:=} \{\tilde{X}_t =X_{T-t} ; 0 \leq t \leq T\}
% \end{equation} An SDE for $ \tilde X$ has been obtained in \cite{haussmann1986time} under the following integrability assumption about the density. Similar conclusions have also been reached in \cite{anderson1982reverse,cattiaux2023time} and \cite{follmer2005entropy}.
% \begin{assumption}\label{asum:bound}
%     For any open bounded set $\Omega \in \Re^n$, the following integrability condition holds
    % \begin{equation}\label{eq:asum-bound}
    %     \int_0^T \int_\Omega |p_t(x)|^2 + \|g(t,x)^\top \partial_x p_t(x)\|^2 \ud x \ud t < \infty ,
    % \end{equation}
%     where $\partial_x p_t(x)$ denotes the partial derivative of $p_t(x)$ with respect to $x$. $g: [0,T]\times \R^n \rightarrow \R^{n \times m}$ is assumed to be smooth and globally Lipschitz. We also assume $G(t,x)\mathrel{:=} g(t,x)g(t,x)^\top  \in L^2(p_t)$ for all $t \in [0,T]$.
% \end{assumption}

% Let $ G(t,x) \mathrel{:=} g(t,x)g(t,x)^\top$. 
% \begin{proposition}\label{prop:tr-sde}
%  Consider the stochastic process $X$ that solves the SDE \eqref{eq:fsde} and its time reversal $\tilde X$ defined according to \eqref{eq:tr-def} under Assumption \ref{asum-fg} and \ref{asum:bound}. Then
%     \leavevmode
%     \begin{enumerate}
%     \item The stochastic process $\tilde X$ admits the SDE
%         \begin{subequations}
%         \begin{align}
%             \ud \tilde{X}_t = (&-h(T-t,\tilde{X}_t) +\epsilon^2 b(T-t,\tilde{X}_t))\ud t \nonumber
%             \\&+ \epsilon g(T-t,\tilde{X}_t) {\ud}  \tilde W_t,  \quad \tilde{X}_0 \sim p_T       \label{eq:tr-sde} 
%         \end{align}
%         where 
%         \begin{align}\label{eq:score-def}
%             b(t,x) \mathrel{:=} \frac{1}{p_t(x)} \partial_x(G(t,x)p_t(x)),
%         \end{align}
%         % with $ G(t,x) \mathrel{:=} g(t,x)g(t,x)^\top$. We assume $G(t,x) \in L^2(p_t)$
%     \end{subequations}
%     \item \label{prop:score-eq}$b(t,x)$ is the solution to the minimization problem
%     \begin{align}\label{eq:score-opt}
%         \min_{\psi(\cdot)\in H^1(p_t)} \mathbb{E}\left[ \frac{1}{2}\|\psi(t,x)\|^2 + \trace(G(t,x) \partial_x \psi(t,x)) \right]
%     \end{align}
%     \end{enumerate}
% \end{proposition}

% \begin{proof}
%     \begin{enumerate}
%         \item Under Assumption \ref{asum-fg} and Assumption \ref{asum:bound}, the time-reversal SDE \eqref{eq:tr-sde} and equation \eqref{eq:score-def} have been established in \cite[Thm. 2.1.]{haussmann1986time}. 
%         \item For convenience, we denote $b(t,x,p)$ as $b$, and similarly, we abbreviate $G(t,x)$ and $p_t(x)$ in following derivation.  Let $J(b) = \mathbb{E}\left[ \frac{1}{2}\|b\|^2 + \trace(G \partial_x b) \right]$ and $b^\star = \frac{1}{p_t} \partial_x(Gp_t)$. By applying the integration by parts and following the \cite[Thm. 1]{hyvarinen2005estimation}, we can rewrite the objective function as 
%         % We apply integration by parts to $\Expect [\trace(G \partial_x b)]$. This yield
%         % \begin{align*}
%         %     \Expect[\trace(G \partial_x b)] &= \int \trace(G\partial_x b) p_t \ud x \\
%         %     &= - \int b^\top \partial_x (G p_t) \ud x \\
%         %     &= - \Expect\left[b^\top \frac{1}{p_t} \partial_x (G p_t)\right]
%         % \end{align*}
%         \begin{align}
%             J(b) &= \Expect \left[\frac{1}{2}\|b-b^\star\|^2\right] + C
%         \end{align}
%         where $C=-\frac{1}{2} \Expect[\|b^\star\|^2]$ independent of $b$. Hence, $b^\star$ is the unique minimizer of $J(b)$, concluding the proof.
%         % \begin{align}
%         %     J(b) &= \Expect \left[ \frac{1}{2}\|b\|^2 - b^\top b^\star\right] \nonumber \\
%         %     &= \Expect \left[\frac{1}{2}\|b-b^\star\|^2 - \frac{1}{2}\|b^\star\|^2\right] \nonumber \\
%         %     & \geq -\frac{1}{2} \Expect[\|b^\star\|^2] \label{eq:jb-ineq}
%         % \end{align}
%     % The equality holds if and only if $b=b^\star$.
%     \end{enumerate}
% \end{proof}

% For any perturbation $c \in C^2$ and scalar $\alpha \in \R$ s.t. $\int |c| < \infty$. 
%         \begin{equation}
%         J(b^\star+\alpha c) \!=\! \Expect \left [\frac{1}{2}\|b^\star\!+\!\alpha c\|^2 \!+\! \trace (G \partial_x( b^\star + \alpha c))   \right] \nonumber
%     \end{equation}
%     \begin{align*}
%         J(b^\star+\alpha c) \!-\! J(b^\star) \!&=\! \alpha( {b^\star}^\top c + \trace(G\partial_x c)) \!+\! \frac{1}{2}\alpha^2 \|c\|^2
%     \end{align*}
%     According to the first-order necessary condition for optimality, we have
%     \begin{equation}\label{eq:1-opt}
%         \Expect \left [ {b^\star}^\top c + \trace(G \partial_x c) \right] = 0 
%     \end{equation}
    
%      Apply integration by parts to $\Expect [\trace(G \partial_x c)]$
%     \begin{align}
%         \Expect [\trace(G \partial_x c)] &= \int \trace(G\partial_x c) p_t(x)\ud x \nonumber \\
%         &= - \int c^\top \partial_x(G p_t(x)) \ud x \nonumber \\
%         &= \Expect [  c^\top( -\frac{1}{p_t(x)}\partial_x(G p_t(x)))] \label{eq:int-part}
%     \end{align}
% Here, we derive $b$ is the solution to \eqref{eq:score-opt} using the calculus of variation method.
%     Define $J(b) \mathrel{:=} \mathbb{E}\left[ \frac{1}{2}\|b(t,x,p)\|^2 + \trace(G(t,x) \partial_x b(t,x,p)) \right]$, Let $b^\star$ be the minimizer of objective function which means $J(b^\star) = \min J(b)$. For any perturbation $c \in C^2$ and scalar $\alpha \in \R$ s.t. $\int |c| < \infty$
%     \begin{equation}
%         J(b^\star+\alpha c) = \Expect \left [\frac{1}{2}\|b+\alpha c\|^2 + \trace (G \partial_x b) +\alpha \trace (G \partial_x c)   \right] \nonumber
%     \end{equation}
%     Considering the first order optimal necessary condition, we can get
%     \begin{equation}\label{eq:1-opt}
%         \Expect \left [\langle b, c\rangle + \trace(G \partial_x c) \right] = 0 
%     \end{equation}
%     Apply integration by parts to $\Expect [\trace(G \partial_x c)]$
%     \begin{align}
%         \Expect [\trace(G \partial_x c)] &= \int \trace(G\partial_x c) p_t(x)\ud x \nonumber \\
%         &= - \int c^\top \partial_x(G p_t(x)) \ud x \nonumber \\
%         &= \Expect [ \langle c, -\frac{1}{p_t(x)}\partial_x(G p_t(x))\rangle] \label{eq:int-part}
%     \end{align}
%     Combining \eqref{eq:1-opt} and \eqref{eq:int-part}, we can conclude our proposition.
% \begin{remark}
%     The SDE \eqref{eq:tr-sde} is the time-reversal formulation of \eqref{eq:fsde}. The optimization \eqref{eq:score-opt} is known as implicit score matching \cite{tang2024score}. $b(t,x)$ is known as F$\Ddot{\text{o}}$llmer's drift term or score function \cite{follmer2005entropy,hyvarinen2005estimation}.
% \end{remark}


% \begin{lemma}
%     $b(t,x,p)$ is the solution to the minimization problem
% \begin{equation}
%     \min_{b(\cdot)} \mathbb{E}\left[ \frac{1}{2}\|b(t,x,p)\|^2 + \trace(G(t,x) \partial_x b(t,x,p)) \right]
% \end{equation}
% \end{lemma}
% \begin{proof}
    
% \end{proof}

% Following our definitions of $\cev X$ and $X$, we can conclude that these two stochastic processes have equal probability law, i.e. $X \overset{d}{=} \cev X$.
% Consider two stochastic process $\vec{X} = \{\vec{X_t};0\leq t \leq T\}$ and $\cev{X} = \{\cev{X_t}; 0 \leq t \leq T\}$ are governed by the following forward backward SDE (FBSDE)
% \begin{subequations}
% \begin{align}
%     &\ud \vec{X_t} \!=\! f(t,\vec{X_t})\ud t + g(t,\vec{X_t})  \vec{\ud} W_t,  \vec{X_0} \sim p_0, \\
%     &\ud \cev{X_t} \!=\! f(t,\cev{X_t})\ud t \!-\! b(t,\cev{X_t},p_{X_t})\ud t + g(t,\cev{X_t}) \cev{\ud} W_t,  \cev{X_T} \sim p_T
% \end{align}  
% where,
% \begin{align}
%     b(t,x,p) \mathrel{:=} \frac{1}{p_t(x)} \partial_x(G(t,x)p_t(x)),
% \end{align}
% \end{subequations}
% $ G(t,x) \mathrel{:=} g(t,x)g(t,x)^\top$, $p_t(x)$ is the probability density function of $X_t$, and $\vec{X_t} \overset{d}{=} \cev{X_t}$. $b(x,p)$ is known as F$\Ddot{\text{o}}$ller's drift term or score function. Moreover $b(t,x,p)$ is the solution to the minimization problem
% \begin{equation}
%     \min_{b(\cdot)} \mathbb{E}\left[ \frac{1}{2}\|b(t,x,p)\|^2 + \trace(G(t,x) \partial_x b(t,x,p)) \right]
% \end{equation}

% Assume $k(t,X_t)$ is the solution to Problem \ref{prob:p1}, which means system state evolves forward according to the forward SDE 
% \begin{equation}\label{eq:fsde}
%     \ud \vec{X_t} = f(t,\vec{X_t})\ud t + g(t,\vec{X_t}) (k(t,\vec{X_t}) \ud t + \epsilon \vec{\ud} W_t),
% \end{equation}
% and it satisfies $\vec{X_0}=x$, $\vec{X_T} = y$.

% \at{In the first subsection, present the time-reversal theory, without any control input. In the next subsection, present the method}
\subsection{Solution to problem~\ref{prob:steer-to-y}}\label{sec:naive-alg}
% \at{We use a time-reversal methodology to solve Problem~\ref{prob:p1}. In order to do so, we introduce a stochastic process $Z_t$ that is initialized at the target $y$ and follows the SDE
% \begin{align}\label{eq:alg-tr-sde}
%     \ud Z_t = &\tilde f(t,Z_t)\ud t +  \epsilon\tilde g(t,Z_t) \ud W_t,~~ \tilde X_0=y
% \end{align}
% where the the functions $\tilde f$ and $\tilde g$ will be designed such that the time-reversal of $Z_t$ solves problem~\ref{prob:p1}. According to some result the time-reversal $\tidle Z =??$ satisfies the SDE.... In order to match this 
% }


We propose to solve Problem~\ref{prob:steer-to-y} using the time-reversal theory that is presented in Section~\ref{sec:tr}. We start with the process $Z_t$ from~\eqref{eq:fsde} and initialize it  at $Z_0=x_f$. With this initialization, the reversed process $\tilde Z_t:=Z_{T-t}$ satisfies  the terminal condition $\tilde Z_T=Z_0=x_f$, and its dynamics has an absorption property for $x_f$. Therefore, in order to solve Problem~\ref{prob:steer-to-y}, we design $h$ so that SDE~\eqref{eq:tr-sde} for $\tilde Z_t$ takes a form similar to~\eqref{eq:sde} for $X_t$. Using the decomposition of the score function as
\begin{equation}\label{eq:score-decom}
    s(t,x) =  g(x) k^\star(t,x) + \mathfrak g(x)
\end{equation}
where the $i$-th component of $ k^\star:[0,T] \times \R^n \rightarrow \R^m$ and  $\mathfrak g :  \R^n \rightarrow \R^n$ are defined as 
\begin{align}\label{eq:k-star}
     k_i^\star(t,x) &\mathrel{:=} \frac{1}{ p(t,x)} \sum_{j=1}^n \partial_{x_j}(g_{j,i}(x)  p(t,x)),\\
 \mathfrak g_i(t,x) &\mathrel{:=} \sum_{j=1}^n \sum_{k=1}^m g_{j,k}(x) \partial_{x_j}g_{i,k}(x), \label{eq:frak-g}  
\end{align}
% and the $i$-th component of $\mathfrak g :  \R^n \rightarrow \R^n$ is defined as 
% \begin{equation}
%     \mathfrak g_i(t,x) \mathrel{:=} \sum_{j=1}^n \sum_{k=1}^m g_{j,k}(x) \partial_{x_j}g_{i,k}(x). 
% \end{equation}
the SDE~\eqref{eq:tr-sde} takes the form 
\begin{align*}
    \ud \tilde Z_t = 
    &\big(-h(\tilde Z_t) + \epsilon^2\mathfrak g(\tilde Z_t)\big)\ud t \\&+  g(\tilde Z_t) \big(\epsilon^2 k^\star(T-t,\tilde Z_t)\ud t + \epsilon \ud \tilde W_t\big).
\end{align*}
Designing the function $h$ and control law $k$ according to 
\begin{align}\label{eq:tilde-f}
    h(x) &:= -f(x) +  \epsilon^2 \mathfrak  g(x),\\
    k(t,x) &:= \epsilon^2 k^\star(T-t,x),\label{eq:fbcontrol}
\end{align}
yields the following expressions for SDEs of $Z_t$ and $\tilde Z_t$:
\begin{subequations}
\begin{align}\label{eq:Z-alg}
\ud Z_t &= (-f(Z_t) + \epsilon^2 \mathfrak g(Z_t))\ud t +  \epsilon g(Z_t) \ud W_t,~Z_0=x_f
\\ \ud \tilde Z_t &= 
    f(\tilde Z_t)\ud t+  g(\tilde Z_t) (k(t,\tilde Z_t)\ud t + \epsilon \ud \tilde W_t).\label{eq:tilde-Z-alg}
\end{align}
Moreover, using the control law~\eqref{eq:fbcontrol} in~\eqref{eq:sde} concludes the following SDE for $X_t$: 
\begin{align}\label{eq:X-alg}
    \ud X_t = f(X_t)\ud t\!+ \! g(\tilde X_t) (k(t, X_t)\ud t + \epsilon \ud W_t),\quad X_0=x_0.
\end{align}
\end{subequations}
In summary, we have constructed three stochastic processes: 
\begin{enumerate}
    \item The process $Z_t$ that solves~\eqref{eq:Z-alg} from initial condition $Z_0=x_f$. The density of this process is denoted by $p(t,x)$;
    \item The process $\tilde Z_t := Z_{T-t}$ that solves~\eqref{eq:tilde-Z-alg} and satisfies the condition $\tilde Z_T=Z_0=x_f$. The control law $k(t,x)$ is defined in~\eqref{eq:fbcontrol} and~\eqref{eq:k-star};
    \item The process $X_t$ that solves~\eqref{eq:X-alg} starting from the initial condition $X_0=x_0$.   
\end{enumerate}
Note that, the SDE for $\tilde Z_t$ has the same form as the SDE for $X_t$ and the control law steers the process $\tilde Z_t$ to the target state $\tilde Z_T=x_f$. However, it remains to be shown that the control law steers $X_t$ to $X_T=x_f$, despite the difference in the initial conditions; $\tilde Z_0 = Z_T$ is random with probability density function $p(T,x)$, whereas $X_0=x_0$ is deterministic. The next result shows that $X_t$ reaches the target state $X_T=x_f$  whenever the initial condition satisfies $p(T,x_0)>0$.  

% for our control problem with the feedback control law .  The terminal condition $\tilde Z_T=x_f$ is satisfied by the time-reversal relationship $\tilde Z_t=Z_{T-t}$ and initializing $Z_0=x_f$. 
% we design the function $h$ such that 
% \begin{align}
%     -h(x) + \epsilon^2 s(T-t,x) &= f(x) + g(x)k(t,x)\\
%      -h(x) + \epsilon^2 s(T-t,x) &= f(x) + g(x)k(t,x)
% \end{align}
% for some control law $k:[0,T]\times \R^n \to \R^m$. 


% Therefore, the drift function $\tilde h$ in~\eqref{eq:tr-sde} serves as a ``control law'' that steers the process $\tilde Z_t$ to the target state $x_f$. We use this idea to solve Problem~\ref{prob:steer-to-y}. 

% We design the drift function $h$ such that the drift function $\tilde h$ of the reversed process~\eqref{eq:tr-sde} takes the form $f+gk$ for some control law $k$. In particular, we want to deisgn

% a process $\tilde X$ that is initialized at the target $y$ and follows the SDE
% \begin{align}\label{eq:alg-tr-sde}
%     \ud \tilde X_t = &\tilde f(\tilde X_t)\ud t +  \epsilon g(\tilde X_t) \ud \tilde W_t,~~ \tilde X_0=y
% \end{align}
% The drift function $\tilde f$ is  is designed such that the reversed process $X_t:=\tilde X_{T-t}$ solves problem~\ref{prob:steer-to-y}. In particular, the terminal condition is immediately satisfied because $X_T=\tilde X_0=y$. Next, we check if $X_t$ satisfies SDE~\eqref{eq:sde} for some feedback control law $U_t=k(t,X_t)$. If the reversibility conditions, Assumption~\ref{asum-fg} and \ref{asum-brackets}, are valid when $h$ is replaced by $\tilde f$, the reversed process $X_t:=\tilde X_{T-t}$ satisfies the SDE
% \begin{align}
%     \ud X_t = (-\tilde f(X_t)+\epsilon^2 s(T-t,X_t)) \ud t + \epsilon g(X_t) \ud W_t,  \label{eq:alg-fsde}
% \end{align}
% where  $s:(0,T] \times \Re^n \to \Re^n$ is the score function whose $i$-th component is given by
%   \begin{align}\label{eq:score-def}
%             s_i(t,x) \mathrel{:=} \frac{1}{\tilde p(t,x)} \sum_{j=1}^n  \partial_{x_j}(G_{i,j}(x)\tilde p(t,x)).
%             %\quad %\forall (t,x)\in (0,T]\times \R^n 
%         \end{align}
% and $\tilde p(t,x)$ denotes the probability density function of $\tilde X_t$. 
% \at{Explain why this design would solve problem 1.}
% where $\tilde p_0$ denote the probability density of the terminal time system state $\tilde X_T$ generated by SDE \eqref{eq:alg-tr-sde}. 
% , we know that forward process $X$ admits the following SDE
% \begin{align}
%     \ud X_t = (-\tilde f(t,X_t)+b(t,X_t)) \ud t + \epsilon g(t,X_t) \ud W_t, \nonumber\\~~X_0\sim \tilde{p}_T, \label{eq:alg-fsde}
% \end{align}
% where
% \begin{equation}
%     b(t,x) = \frac{\epsilon^2}{\tilde p_{T-t}(x)}\partial_x(G(t,x)\tilde p_{T-t}(x))
% \end{equation}
% According to the definition of the time-reversal process $X_t = \tilde X_{T-t}$, so we have $X_T = y$.
% The function $\tilde f$ is designed such that  the forward process $X$ satisfies the SDE
% \begin{align}\label{eq:target-sde}
%     \ud X_t = f(t,X_t)\ud t + g(t,X_t) (&k(t,X_t) \ud t + \epsilon \ud W_t)
% \end{align}
% for some control law $k:[0,T] \times \R^n \rightarrow \R^m$. 
% \at{Here you should use (10) and (11) to conclude a relationship for $\tilde f$ and $k$. Then, to obtain $\tilde f$ and $k$, you decompose $b$ (13). You use the result ton conclude the formula for $k$ and $\tilde f$ (17).}
% Comparing SDE~\eqref{eq:alg-fsde} with \eqref{eq:sde}, where $U_t$ is relaced by its control law $k(t,X_t)$, yields the following relationship: 
% \begin{equation}\label{eq:design-eq}
%     f(x)+g(x)k(t,x) = -\tilde f(x) + \epsilon^2 s(T-t,x)
% \end{equation}
% The relationship is used to obtain the drift function $\tilde f$ used in~\eqref{eq:alg-tr-sde} and the control law $k$ used to solve problem~\ref{prob:steer-to-y}. In order to do so, we decompose the score function $s(t,x)$ into two parts as 
% \begin{equation}
%     s(t,x) =  g(x) k^\star(t,x) + \mathfrak g(x)
% \end{equation}
% where the $i$-th component of $ k^\star:[0,T] \times \R^n \rightarrow \R^m$ is defined as 
% \begin{equation}
%      k_i^\star(t,x) \mathrel{:=} \frac{1}{\tilde p(t,x)} \sum_{j=1}^n \partial_{x_j}(g_{j,i}(x) \tilde p(t,x))
% \end{equation}
% and the $i$-th component of $\mathfrak g :  \R^n \rightarrow \R^n$ is defined as 
% \begin{equation}
%     \mathfrak g_i(t,x) \mathrel{:=} \sum_{j=1}^n \sum_{k=1}^m g_{j,k}(x) \partial_{x_j}g_{i,k}(x). 
% \end{equation}
% % Letting $-\tilde f(t,X_t) + b(t,X_t) = f(t,X_t)+g(t,X_t)k(t,X_t)$, we obtain the equation
% % \begin{equation}
% %     f+gk  = \epsilon^2 g k^\star + \epsilon^2 \tilde g - \tilde f
% % \end{equation}
% The decomposition, along with the relationship \eqref{eq:design-eq}, imply that 
% \begin{subequations}
%     \begin{align}\label{eq:tilde-f}
%         \tilde f(x) &= -f(x) + \epsilon^2 \mathfrak  g(x)\\
%         k(t,x) &= \epsilon^2 k^\star(T-t,x)\label{eq:fbcontrol}
%     \end{align}
% \end{subequations}
% With these definitions, the SDE~\eqref{eq:alg-fsde} becomes 
% \begin{align}
%     \ud X_t = f(X_t)+ g(X_t)(k(t,X_t)\ud t + \epsilon \ud W_t). 
% \end{align}
% and the terminal condition $X_T=y$ is satisfied by the time-reversal relationship $X_T=\tilde X_0=y$. 
% The probability density of the initial state $X_0=\tilde X_t \sim \tilde p(T,x)$. The next result shows that the terminal condition $X_T=y$ is satisfied as long as $X_0$ is initialized from a point $x_0$ with non-zero probability density, i.e. $\tilde p(T,x_0)>0$. 


% As a result, if the assumptions for Proposition \ref{prop:tr-sde} are valid, we constructed a control law such that if $X_0 \sim \tilde p_T$ then $X_T = y$. In the following proposition, we show that this is true if $X_0$ starts from any initial condition $x$ such that $P(\tilde X_T = x)>0$, then $X_T=y$, implying $U_t = k(t, X_t)$ is the solution to Problem \ref{prob:p1}.
% \at{Before the proposition, state the conclusion you are getting from this procedure. The fact that if $X_0\sim \tilde p_0$, then $X_T=y$. This follows from Prop. 1. Then, discuss that in the proposition, you want to draw conclusion if $X_0$ is initilized at $x$.}
\begin{proposition}\label{prop:sol2p1}
% Consider problem~\ref{prob:steer-to-y} for the stochastic control system~\eqref{eq:sde}. 
Let $p(t,x)$ denote the probability density function  of $Z_t$ defined according to~\eqref{eq:Z-alg} and define the control law $k(t,x)$ according to~\eqref{eq:fbcontrol} and~\eqref{eq:k-star}. If the initial condition $X_0=x_0$ satisfies $p(T,x_0)>0$, then, problem~\ref{prob:steer-to-y} is solved with the feedback control law $k(t,x)$. 
    % The control input $U_t = k(t,X_t)$ can bring any initial state $\tilde x$, s.t. $P (\tilde X_T=\tilde x) >0$, to target state $y$ which means $U_t$ is the solution to Problem \ref{prob:p1}.
    % Assuming the conditions of Proposition \ref{prop:tr-sde} are valid for SDE \eqref{eq:alg-tr-sde} and probability density $\tilde p_t$, then Problem 1 is solved with the feedback control law \eqref{eq:fbcontrol}, if $X_0=x$ with $P(\tilde X_T = x)>0$.   
\end{proposition}
\begin{proof}
% Let $P_{\tilde X_{t},\tilde X_T}(x,x')$ denote the join probability density function for $(\tilde X_{t},\tilde X_T)$ where $\tidle X_t$ is defined according to SDE~\eqref{eq:alg-tr-sde}.
Let $\kappa_{t,s}(x'|x)$ and $\tilde \kappa_{t,s}(x'|x)$ be the probability transition kernels for SDEs~\eqref{eq:Z-alg} and~\eqref{eq:tilde-Z-alg}, respectively, as defined in Section~\ref{sec:kernel}.  
The probability transition kernel associated with the SDE~\eqref{eq:X-alg} is also $\tilde \kappa_{t,s}(x'|x)$ due to the fact that SDEs~\eqref{eq:tilde-Z-alg} and~\eqref{eq:X-alg} are identical.  Therefore, with the initial condition $X_0=x_0$, the probability density function of $X_t$ becomes equal to $\tilde \kappa_{t,0}(x|x_0)$. The goal is to show that  $\tilde \kappa_{t,0}(x|x_0)$ approaches the Dirac delta distribution $\delta_{x_f}(x)$ (in the weak sense) as $t$ approaches $T$. The identity~\eqref{eq:kappa-identity} implies  
% \begin{align*}
%     \tilde \kappa_{t,0}(x|x_0)\tilde p(t,x) &= \kappa_{T-t,0}(x|x_0) p(0,x_0) \\ &=\kappa_{T-t,0}(x|x_0) \tilde \kappa_{T,0}(x_0|y)
% \end{align*}
% As a result, if $\kappa_{T,0}(x_0|y)>0$, we have 
    \begin{align*}
   \tilde \kappa_{t,0}(x|x_0)  &= \frac{  \kappa_{T,T-t}(x
    _0|x)p(T-t,x)}{ p(T,x_0)}\\
    &=\frac{  \kappa_{T,T-t}(x
    _0|x)p(T-t,x)}{ \kappa_{T,0}(x_0|x_f)}
\end{align*}
where we used the assumption that $p(T,x_0)>0$ and the fact that $ p(T,x)= \kappa_{T,0}(x|x_f)$ due the initial condition $Z_0=x_f$. 
 Taking the limit as $t\to T$ and using the fact that $p(T-t,x)$ approaches $\delta_{x_f}(x)$ concludes the result. 
\end{proof}
% \subsection{Slack Terminal Constraints}

\subsection{Avoiding singularity of the control law}\label{sec:singular}
The feedback control law~\eqref{eq:fbcontrol} becomes singular in the limit as $t$ approaches the terminal time $T$ because the distribution $p(T-t,x)$ approaches the Dirac delta distribution $\delta_{x_f}(x)$. The singularity is unavoidable when an almost sure constraint $X_T=x_f$ is required. In order to avoid the singularity, we consider Problem~\ref{prob:approx-steer-to-y} where the almost sure equality is relaxed to a bound on the average distance to the target. We propose to solve problem~\ref{prob:approx-steer-to-y} using the time-reversal procedure presented in Section~\ref{sec:naive-alg}, with the difference that the process $Z_0$ is now initialized from a Gaussian distribution around $x_f$, i.e. $\mathcal N(\cdot;x_f,\sigma^2)$ with $\sigma>0$. 
We prove that, with a small enough $\sigma$, the resulting control law is able to solve problem~\ref{prob:approx-steer-to-y} while maintaining non-singularity, in the linear Gaussian setting,  with error bounds established in Proposition~\ref{prop:terminal-error},
% We establish error 
% We prove that for linear systems this approach steers the original process $X$ to within a bounded neighborhood of $x_f$, with error bounds established in Proposition~\ref{prop:terminal-error}, 
while deferring the analysis of error bounds for the general nonlinear case to future work. 


 % We prove this conjecture in the linear Gaussian setting in the next section and leave the analysis of nonlinear setting to future work.  
% \Ali{Given that no actual conjecture is made, I suggest changing the last sentence to: ``We prove that for linear systems this approach steers the original process $X$ to within a bounded neighborhood of $x_f$, with error bounds established in Proposition~\ref{prop:terminal-error}, while deferring the analysis of error bounds for the general nonlinear case to future work.''.}

\subsection{Analysis of the linear Gaussian setting}\label{sec:lin-gau}
% \at{Outline:
% \begin{enumerate}
%     \item The first goal is to find the expression for the control law in the linear setting. State the SDE for $Z_t$. Find the kernel. Find the density. Use it to express the control law $k$. Define what you mean by $N(x;m\Sigma)$
%     \item Show that the control law becomes singular as $t\to T$. In order to resolve this issue, you start from a Gaussian distribution around $x_f$. State the new density. 
%     \item Use (15) to compute the kernel and density for $X_t$. 
% \end{enumerate}
% }
In this section, we study the proposed time-reversal method in the linear Gaussian setting, where
\begin{align*}
    f(x) = Ax,\quad g(x) = B,
\end{align*}
for matrices  $A \in \R^{n\times n}$ and $B \in \R^{n \times m}$. In this case, $f$ is linear and $g$ is constant, implying that both functions are smooth and globally Lipschitz. Therefore, Assumption~\ref{asum-fg} is satisfied. Furthermore, Assumption \ref{asum-brackets} will also hold whenever $(A,B)$ is controllable.
% \at{Talk about the fact that f and g satisfy assumption 1. And they will satisfy assumption 2 when A B is controllable.}


Under such setting, the SDEs \eqref{eq:Z-alg}, \eqref{eq:tilde-Z-alg}, and ~\eqref{eq:X-alg} take the form
\begin{subequations}
\begin{align}
    &\ud Z_t = -A Z_t \ud t + \epsilon B \ud W_t, \hspace{59pt} Z_0 = x_f,\label{eq:Z-lin}\\
    &\ud \tilde Z_t = A \tilde Z_t \ud t + B( k(t,\tilde Z_t) \ud t +\epsilon \ud \tilde W_t), \hspace{2pt} \tilde Z_T = x_f, \label{eq:tilde-Z-lin}\\
    &\ud X_t \!=\! A X_t \ud t \!+\! B(k(t,X_t)\ud t \!+ \!\epsilon \ud W_t), \hspace{10pt} X_0\! =\! x_0. \label{eq:X-lin}
\end{align}
\end{subequations}
% where $A \in \R^{n\times n}$ and $B \in \R^{n \times m}$. 
% Following the same procedure as in the proof of Proposition \ref{prop:sol2p1}, Let $\kappa_{t,s}(x^\prime|x)$ denote the probability transition kernel from time $s$ to time $t$ associated with SDE \eqref{eq:Z-lin} and $\tilde \kappa_{t,s}(x^\prime|x)$ denote the probability transition kernel associated with SDE \eqref{eq:tilde-Z-lin}. We can obtain the distribution of $\tilde Z_T$ given the initial state $\tilde Z_0 = x_0$
The probability density of $Z_t$ is Gaussian, for all $t>0$, because \eqref{eq:Z-lin} is linear and the initial state $Z_0$ is deterministic.  In particular, $p(t,x)=\mathcal N(x;m_t,\Sigma_t)$ where the mean and covariance are given by
\begin{align*}
      m_t &= e^{-At}x_f,\\
      \Sigma_t &=  \epsilon^2 \int_0^t e^{-A(t-s)}BB^\top e^{-A^\top(t-s)}\ud s.
\end{align*}
Substitution of the Gaussian distribution formula in \eqref{eq:k-star}, and the use of \eqref{eq:fbcontrol}, yields the following formula for the feedback control law:
\begin{align}\label{eq:k-linear}
    k(t,x) 
    &= -\epsilon^2 B^\top \Sigma_{T-t}^{-1}(x-m_{T-t}),
\end{align}
for $t \in [0,T)$. This is similar to the control law that appears in~\cite[Eq. (11)]{pakniyat2021steering} and \cite{chen2015stochastic}, obtained through a stochastic optimal control formulation.
% \at{Add refereecnes to Ali's and Yongxins paper. Point to equation numbers}

It is worth remarking that in the limit as $t \to T$, the covariance $\Sigma_{T-t} \to 0$, thus the control law becomes singular, as described in Section~\ref{sec:singular}. To resolve the singularity issue, we initialize $Z_0 \sim \mathcal{N}(\cdot;x_f,\sigma^2I)$. Under this initial condition, the distribution of $Z_t$ remains Gaussian $\mathcal N(x;m_t,Q_t)$ with the same mean  as before, but with a new covariance that is given by 
$
    Q_t = \Sigma_t + \sigma^2 e^{-At}e^{-A^\top t}
$.
The resulting feedback control law is
\begin{align}\label{eq:approx-k-linear}
    \hat k(t,x) 
    &= -\epsilon^2 B^\top Q_{T-t}^{-1}(x-m_{T-t}).
\end{align}
The new feedback control law remains nonsingular as $t\to T$. However, there is no guarantee that it would steer $X_t$ to the target state $x_f$. The following proposition  characterizes the error $\mathbb E[\|X_T-x_f\|^2]$ when the control law~\eqref{eq:approx-k-linear} is used instead of~\eqref{eq:k-linear}. 
\begin{proposition}\label{prop:terminal-error}
    % \at{Express the result (16) as a proposition and re-write the derivations as the proof. Express the fact that in the limit $\sigma \to 0$, the error goes to zero, as the second part of the proposition.}
    Under the feedback control law defined in \eqref{eq:approx-k-linear}, the expected squared error between the terminal state and the target state is given by
    \begin{equation}\label{eq:error}
        \Expect[\|X_T-x_f\|^2] = \sigma^4\|e^{TA}x_0-x_f\|^2_{M^{-2}} + \sigma^2(n -\trace(M^{-1})),
    \end{equation}
    where $M := \sigma^2 I +  \epsilon^2\int_0^T e^{As}BB^\top e^{A^\top s}\ud s$. Moreover, for any $\delta>0$, there exists a small enough $\sigma>0$ that solves Problem~\ref{prob:approx-steer-to-y}.
\end{proposition}
\begin{proof}
     Let $\kappa$ and $\tilde \kappa$ denote the probability transition kernels associated with SDEs~\eqref{eq:Z-lin} and~\eqref{eq:tilde-Z-lin}, respectively, similar to the proof of Proposition~\ref{prop:sol2p1}. In terms of the kernels, the probability distribution of $X_T$ is equal to $\tilde \kappa_{T,0}(\cdot|x_0)$, because the SDE~\eqref{eq:X-lin} and~\eqref{eq:tilde-Z-lin} have the same form and $X_0=x_0$.  Then, upon the application of the time-reversal relationship~\eqref{eq:kappa-identity}, and the fact that $p(t,x) = \mathcal N(x;m_t,Q_t)$,
\begin{align*}
    \tilde \kappa_{T,0}(x|x_0) &=  \frac{ \kappa_{T,0}(x_0|x)p(0,x)}{p(T,x_0)} \\
    & = \frac{\mathcal N(x_0;e^{-AT}x,\Sigma_T) \mathcal N(x;x_f,\sigma^2I)}{\mathcal N(x_0;m_T,Q_T)}\\&=\mathcal N(x;\mu,P)
\end{align*}
where 
\begin{align*}
    \mu &= x_f + \sigma^2 M^{-1}e^{TA}(x_0 - e^{-TA}x_f),\\
    P & =  \sigma^2(I- \sigma^2 M^{-1}).
\end{align*}
  Now, to compute the error   $\mathbb E[\|X_T-x_f\|^2]$, we use the fact that $X_T \sim \mathcal N(\cdot;\mu, P)$ and, hence: 
\begin{align}
    \mathbb E[\|X_T-x_f\|^2] &= \|\mu-x_f\|^2 + \trace(P) \nonumber
    % \\
    % &= \sigma^4\|e^{TA}x_0\!-\!x_f\|^2_{M^{-2}} \!+\! \sigma^2(n\! -\!\trace(M^{-1})), \nonumber
\end{align}
which yields~\eqref{eq:error}. Moreover, in the limit as $\sigma \to 0$, $M$~converges to the controllability grammarian matrix which is positive definite  under the controllability assumption.  Therefore, taking the limit of~\eqref{eq:error} as $\sigma \to 0$, concludes that the error converges to zero, implying that for any $\delta>0$, there exists a $\sigma>0$ such that the error is smaller than $\delta$. 
\end{proof}
% Consider the system \eqref{eq:lingau-bsde} is linear and the initial distribution is Gaussian, the distribution of $\{\tilde X_t:0\leq t \leq T\}$ is always Gaussian. To obtain the $\tilde p(t,x)$, we only need to find the mean and covariance of $\tilde X_t$. Taking expectations on both sides of \eqref{eq:lingau-bsde}
% \begin{equation}
%     \ud \Expect [\tilde X_t] = -A \Expect[\tilde X_t]\ud t, ~~ \Expect[\tilde X_0] = x_f
% \end{equation}
% Let $m_t$ denote the mean of $\tilde X_t$, we can obtain 
% \begin{equation}
%     m_t = e^{-At}x_f
% \end{equation}
% Let $\Sigma_t$ denote the covariance of $\tilde X_t$, and $\Sigma_t = \Expect [\tilde X_t \tilde X_t^\top] - \Expect[\tilde X_t]\Expect[\tilde X_t]^\top$.
% \begin{align}
%     \ud \Sigma_t = \Expect[\ud \tilde X_t \tilde X_t^\top + \tilde X_t \ud \tilde X_t^\top + \ud \tilde X_t \ud \tilde X_t^\top] - \nonumber \\
%     \ud \Expect[\tilde X_t] \Expect[\tilde X_t]^\top - \Expect[\tilde X_t] \ud \Expect[\tilde X_t]^\top
% \end{align}
% \begin{align}
%     \frac{\ud}{\ud t} \Sigma_t = -A \Sigma_t - \Sigma_t A^\top + \epsilon^2 B B^\top, ~~ \Sigma_0 = \sigma^2I
% \end{align}
% The solution for $\Sigma_t$ is
% \begin{equation}
%     \Sigma_t = \sigma^2 e^{-At}e^{-A^\top t} + \epsilon^2 \int_0^t e^{-A(t-s)}BB^\top e^{-A^\top(t-s)}\ud s
% \end{equation}

% In order to do so, let $\kappa$ and $\tilde \kappa$ denote the probability transition kernels associated with SDEs~\eqref{eq:Z-lin} and~\eqref{eq:tilde-Z-lin}, respectively, similar to the proof of Proposition~\ref{prop:sol2p1}. In terms of the kernels, the probability distribution of $X_T$ is equal to $\tilde \kappa_{T,0}(\cdot|x_0)$, because the SDE~\eqref{eq:X-lin} and~\eqref{eq:tilde-Z-lin} have the same form and $X_0=x_0$.  Then, upon the application of the time-reversal relationship~\eqref{eq:kappa-identity}, and the fact that $p(t,x) = \mathcal N(x;m_t,Q_t)$, we have 
% \begin{align*}
%     \tilde \kappa_{T,0}(x|x_0) &=  \frac{ \kappa_{T,0}(x_0|x)p(0,x)}{p(T,x_0)} \\
%     & = \frac{\mathcal N(x_0;e^{-AT}x,\Sigma_T) \mathcal N(x;x_f,\sigma^2I)}{\mathcal N(x_0;m_T,Q_T)}\\&=\mathcal N(x;\mu,P)
% \end{align*}
% where 
% \begin{align*}
%     \mu &= x_f + \sigma^2 M^{-1}(e^{TA}x_0 - x_f),\\
%     P & =  \sigma^2(I- \sigma^2 M^{-1}),
% \end{align*}
% and $M = \sigma^2 I +  \epsilon^2\int_0^T e^{As}BB^\top e^{A^\top s}\ud s$. Note that, in the limit as $\sigma^2$, the variance of the initial condition $Z_0$, goes to zero, we have $\mu \to x_f$, and $P \to 0$, that is $X_T \to x_f$, as predicted by Proposition~\ref{prop:sol2p1}.  Now, we compute the error   $\mathbb E[\|X_T-x_f\|^2]$ using the fact that $X_T \sim \mathcal N(\cdot;\mu, P)$: 
% \begin{align}
%     \mathbb E[\|X_T-x_f\|^2] &= \|\mu-x_f\|^2 + \trace(P) \nonumber \\
%     &= \|e^{TA}x_0-x_f\|^2_{M^{-2}} + \sigma^2(n -\trace(M^{-1})) \label{eq:error}
% \end{align}
% \at{add the definition of $\|\cdot\|_P$ for any positive definite matrix $P$.}

% In order to find the distribution of $\tilde Z_T$, let $ \kappa _{t,0}(x|x^\prime)$ denote the probability transition kernel form time $0$ to time $t$ associated with SDE \eqref{eq:Z-lin}.
% \begin{equation}
%      \kappa _{t,0}(x|x^\prime) = P( Z_t=x| Z_0=x^\prime),
% \end{equation}
% % where $m_t^\prime \!= \!e^{-At}x^\prime$ and
% % \[
% % N_t \!=\! \epsilon^2 \int_0^t e^{\!-\!A(t-s)}BB^\top e^{\!-\!A^\top(t-s)}\ud s.
% % \]
% Let $\tilde \kappa_{t,0}(x|x_0)$ denote the probability transition matrix kernel from time 0 to time $t$ associated with SDE \eqref{eq:tilde-Z-lin}. According to \eqref{eq:kappa-identity}, we can conclude
% \begin{align}
%     \tilde \kappa_{t,0}(x|x_0) &=  \frac{ \kappa_{T,T-t}(x_0|x)p(T-t,x)}{p(T,x_0)}
% \end{align}
% , the distribution of $X_T$ given $X_0=x_0$ can be expressed as
% \begin{equation}
%    P(X_T=x|X_0=x_0) =  \tilde \kappa_{T,0}(x|x_0) = \mathcal{N}(x;\mu,\Sigma)
% \end{equation}
% Let $M = \epsilon^2\int_0^T e^{As}BB^\top e^{A^\top s}\ud s$
% \begin{equation}
%     \Sigma= \sigma^2(I-\sigma^2(\sigma^2I + M)^{-1})
% \end{equation}
% % \begin{equation}
% %     \Sigma= \left( e^{-A^\top T}N_T^{-1} e^{-AT} + \frac{1}{\sigma^2}I   \right)^{-1}
% % \end{equation}
% and
% \begin{align}
%     \mu = x_f + \sigma^2(e^{A^\top T}M^{-1}x_0 - \sigma^2(M+\sigma^2 I)^{-1}e^{A^\top T}M^{-1}x_0 \nonumber\\- (M+\sigma^2 I)^{-1}x_f)
% \end{align}
% The error for mean is 
% \begin{align}
%     \text{error}_{\mu} = \sigma^2(e^{A^\top T}M^{-1}x_0 - \sigma^2(M+\sigma^2 I)^{-1}e^{A^\top T}M^{-1}x_0 \nonumber\\- (M+\sigma^2 I)^{-1}x_f)
% \end{align}
% \begin{equation}
%     \mu = \Sigma(N_T^{-1}e^{-AT}x_0 + \frac{1}{\sigma^2}x_f)
% \end{equation}
% \at{Compute the error. Use matrix inversion lemma}
% Now we can write the control input as 
% \begin{equation}
%     U_t = -\epsilon^2 B^\top \Sigma_{T-t}^{-1} (X_t-m_{T-t})
% \end{equation}
% The system \eqref{eq:lingau-fsde} is linear and initial condition is a deterministic point, so the distribution of $\{X_t;0<t\leq T\}$ is always Gaussian. Following the same procedure as above, 
% \begin{align}
%     \frac{\ud}{\ud t} \Expect [X_t] \!=\! (A \!-\!\epsilon^2BB^\top \Sigma_{T-t}^{-1})\Expect[X_t] &+ \epsilon^2BB^\top\Sigma_{T-t}^{-1}m_{T-t},\nonumber \\ &\Expect[X_0] = x_0
% \end{align}
% By using the method presented in Section \ref{sec:naive-alg}, we can obtain a feedback control law $U_t = k(t,X_t)$. With the initial condition $X_0 = x$, we want to derive the dynamics of the expectation and variance of $X_t$
% \begin{equation}
%     \Expect [\ud X_t] = \Expect [AX_t \ud t + BU_t \ud t + \epsilon B \ud W_t]
% \end{equation}
% \begin{equation}
%     \frac{\ud}{\ud t} \Expect [X_t] = A \Expect[X_t] + B\Expect[U_t], ~~ \Expect[X_0] = x
% \end{equation}
% The solution for this ODE is 
% \begin{equation}
%     \Expect[X_T] = e^{TA}x + \int_0^T e^{(T-\tau)A}B \Expect[U_\tau] \ud \tau
% \end{equation}
% Let $\Sigma_t$ denote the variance of $X_t$, and $\Sigma_t = \Expect [X_t X_t^\top] - \Expect[X_t]\Expect[X_t]^\top$



\begin{remark}\label{rem:improve-alg}
    The error \eqref{eq:error} comprises both a bias term and a variance term. The variance term is independent of $x_0$ and $x_f$, and bounded by $\sigma^2n$. In contrast, the bias term can be significant when $e^{TA}x_0$ and $x_f$ are far apart. This term arises due to the difference between the mean $\mathbb E[Z_T]=e^{-TA}x_f$ and $x_0$. To address this issue, we allow our proposed methodology the flexibility to incorporate a deterministic control input $u_t$ when $Z_t$ is simulated, and modify the control law to  $U_t = k(t,X_t) + \tilde u_t$ with $\tilde u_t = -u_{T-t}$\footnote{Our method can also incorporate a feedback control law in addition to the deterministic input. We do not pursue this in the paper, as computing a feedback law is considerably more challenging than finding a deterministic input.}. The addition of the deterministic input decreases the bias error by bringing the mean of $Z_T$ and $x_0$ closer. Infact, the difference becomes zero in the linear case when $(A,B)$ is controllable.  The details of this modification to the algorithm appears in~\ref{sec:Algorithm}. 
    
    % In particular, the control input $u_t$ is designed to approximately bring $Z_0=x_f$ to $Z_T=x_0$, for the deterministic setting of the model where $\epsilon=0$.  This control may be obtained using various trajectory optimization techniques where additional penalty on the control norm is incorporated~\cite{}. 
    % In addition to decreasing the bias error, the control input $u_t$ serves as an ``importance sampling" mechanism that guides $Z_t$ to be sampled in areas of the domain where the control law is more relevant to the initial condition $x_0$, thus increasing sampling efficiency. Considering the addition of control input $u_t$, the control law that is designed to steer $X_t$ to $x_f$ is modified to  $U_t = k(t,X_t) + \tilde u_t$ with $\tilde u_t = -u_{T-t}$.   
    
    % In particular, we introduce the term $g(Z_t)u_t$ when simulating $Z_t$ in order to steer $Z_T$ close to $x_0$. This modification results 
\end{remark}
% \Ali{I think Section~\ref{sec:improve-alg} is a distraction to the material of this section, and would better be combined with (and appear after) Section~\ref{sec:Algorithm} in a new section ``Computational Considerations'' (or some other similar title), so to leave Section~\ref{sec:results} exclusively for numerical results.}

% \subsection{Improving the algorithm}\label{sec:improve-alg}
% \at{
% \begin{itemize}
%     \item State the objective that you like to reduce the bias error by including a deterministic control input when $Z_t$ is simulated. 
%     \item Give the new equation for Z and tilde Z and X with determinsitc control
%     \item Explain that the covaraince and control laws do not change but the mean cahnges. 
%     \item Explain that if you follow the same steps you arrive at: give the new mean $\mu$. say variance is the same. 
%     \item Choose the input to minimize the error. 
    
% \end{itemize}

% }
%  The error \eqref{eq:error} comprises both a bias term and a variance term. The variance term is independent of $x_0$ and $x_f$, and bounded by $\sigma^2n$. In contrast, the bias term can be significant when $e^{TA}x_0$ and $x_f$ are far apart.  To address this issue, we introduce a deterministic control input $u_t$ when simulating $Z_t$ in order to steer $Z_T$ close to $x_0$. In particular, we modify the SDEs \eqref{eq:X-lin}, \eqref{eq:Z-lin}, and \eqref{eq:tilde-Z-lin} as follows:
%  \begin{subequations}
% \begin{align} 
%     &\ud Z_t = -A Z_t \ud t + Bu_t \ud t + \epsilon B \ud W_t, ~~ Z_0 = x_f,\label{eq:Zu-lin}\\
%     &\ud \tilde Z_t = A \tilde Z_t \ud t + B (k(t,\tilde Z_t) \ud t +\tilde u_{t} \ud t+\epsilon\ud \tilde W_t). \label{eq:tilde-Zu-lin}\\\label{eq:Xu-lin}
%  &\ud X_t = A X_t \ud t \!+ \!B(k(t,X_t)\ud t+ \tilde u_{t}\ud t + \epsilon \ud W_t),  X_0 = x_0,
% \end{align}
%  \end{subequations}
% where $\tilde u_t := - u_{T-t}$.
%     With this modification, the process $Z_t$ still satisfies the time-reversal relationship $Z_t = \tilde Z_{T-t}$. Because $u_t$ is deterministic, the distribution of $Z_t$ remains Gaussian $\mathcal N (\cdot;m_t^u,Q_t)$ with unchanged covariance and new mean given by:
% \begin{equation*}
%     m_t^u = e^{-tA}x_f + \int_0^t e^{(s-t)A}Bu_s \ud s
% \end{equation*}
% Accordingly, the feedback control law \eqref{eq:approx-k-linear} is updated with this new mean $m^u_t$ instead of $m_t$. Following the same procedure as in the proof of Prop~\ref{prop:terminal-error}, we obtain the probability distribution of $X_T$ to be Gaussian $\mathcal N(\cdot;\mu^u,P)$ with the new mean
% \begin{equation*}
%     \mu^u = x_f + M^{-1}(e^{TA}x_0 + \int_0^Te^{(T-t)A}B\tilde u_{t} \ud t - x_f),
% \end{equation*}
% resulting a new bias error:
% \begin{equation*}
%     \|e^{TA}x_0 + \int_0^Te^{(T-t)A}B \tilde u_{t} \ud t - x_f\|^2_{M^{-2}}.
% \end{equation*}
% The controllability of $(A,B)$ implies that there exists a control $\tilde u_t$ that makes the bias error zero. 
% % there exists a control input $\tilde u_t^\star$ that drives the bias error to zero. This control input $\tilde u_t^\star$ corresponds to the solution of a deterministic linear system with a fixed terminal state constraint. 
% % The system's controllability is guaranteed by Assumption \ref{asum-brackets} when $h$ and $g$ are replaced by their linear counterparts $Ax$ and $B$.

% Although, the motivation for including the additional deterministic control is provided for the linear setting, we propose to adopt this strategy to the nonlinear setting. In this general setting, the control input $u_t$ is designed to approximately bring $Z_0=x_f$ to $Z_T=x_0$, for the deterministic setting of the model where $\epsilon=0$.  This control may be obtained using various trajectory optimization techniques where additional penalty on the control norm is incorporated~\cite{}. 
% % The resulting SDE for $Z_t$ is 
% % \begin{align}\label{eq:Z-improved}
% %     \ud Z_t = -f(Z_t)\ud t + \epsilon^2\mathfrak g(Z_t)\ud t  + g(Z_t)(u_t\ud t + \epsilon  g(Z_t)\ud W_t)
% % \end{align}
% In addition to decreasing the bias error, the control input $u_t$ serves as an ``importance sampling" mechanism that guides $Z_t$ to be sampled in areas of the domain where the control law is more relevant to the initial condition $x_0$, thus increasing sampling efficiency. Considering the addition of control input $u_t$, the control law that is designed to steer $X_t$ to $x_f$ is modified to  $U_t = k(t,X_t) + \tilde u_t$ with $\tilde u_t = -u_{T-t}$.   
% The corresponding modified SDEs \eqref{eq:fsde}, \eqref{eq:Z-alg}, and \eqref{eq:tilde-Z-alg} becomes
% \begin{equation}\label{eq:fsde-u}
%     \ud X_t \!=\! f(X_t) \ud t +  g(Z_t)(k(t,X_t) \ud t+u_t \ud t + \epsilon \ud W_t), X_0=x_0,
% \end{equation}
% \begin{align}
% \ud Z_t = (-f(Z_t) + \epsilon^2 \mathfrak g(Z_t))\ud t \!+\! g(Z_t)(&\!-\!u_{t} \ud t + \epsilon  \ud W_t),\nonumber
% \end{align}
% and modify the control input used to simulate $X_t$ to  $U_t = k(t,X_t) + u_t$. The error analysis of this setting is subject to future work. Moreover, $u_t$ improves the exploration efficiency compared to the unforced dynamic. Based on our discussion, we can conclude the algorithm.
\begin{algorithm}[bt]
\caption{Time-Reversal Control Synthesis} 
\begin{algorithmic}[1]
\STATE \textbf{Input:}  sample size $N$, step-size $\Delta t$, variance $\sigma$, deterministic control input $u_t$, function class $\Psi$.   
\STATE  $\{Z_0^i\}_{i=1}^N \sim \mathcal{N}(x_f, \sigma^2I)$

\FOR{$t\in \{\Delta t, 2\Delta t,\ldots,T-\Delta t,T \}$}
\STATE $\{\Delta  W^i_t\}_{i=1}^N\sim N(0,\Delta t I_n)$
	\STATE $ Z^{i}_{t+\Delta t} \!=\!  Z^i_{t}\!+\!(-f(Z^{i}_{t})\!+\!\epsilon^2\mathfrak g( Z^{i}_{t})+g(Z^{i}_{t})u_t)\Delta t  + \epsilon g(Z^{i}_{t})  \Delta  W^i_t$
\ENDFOR
\STATE $k^\star(t,\cdot)\!=\! \argmin_{k\in \Psi} \frac{1}{N}\sum_{i=1}^N \big[\frac{1}{2}\|g( Z^{i}_{t})k(t,Z^i_{t}) \!+\! \mathfrak g( Z^i_{t})\|^2 \!+\!  \sum_{j,l}^n(G_{j,l}(Z^i_{t})\partial_{x_l} (g(Z^{i}_{t})k(t,Z^i_{t})\!+\!\mathfrak g( Z^i_{t}))_j)\big]$

 \STATE \textbf{Output:}  $\{k^\star(t,\cdot)\}_{t\in \{0,\Delta t,\ldots,T\} }$
\end{algorithmic}
\label{alg:improved}
\end{algorithm}





% If $\epsilon=0$, the problem becomes a finite time deterministic control problem.

% \begin{problem}\label{prob:p2}
%     Consider a deterministic control system with the dynamic model 
%     \begin{equation}\label{eq:nonoise_sys}
%         \dot{x}_t = f(t,x_t) + g(t,x_t)u_t, ~~ x_0=x
%     \end{equation}
%     Given $y \in  \R^n$ and a fixed terminal time $T>0$, find a control input $\{u_t;t\in[0,T]\}$ such that  $x_T = y$.
% \end{problem}

% Several methods have been proposed to solve this problem, for example, the trajectory optimization method \cite[Ch. 10]{underactuated} can be applied. Trajectory optimization discretizes the system over the time dimension and uses the nonlinear optimization method to obtain the control input. However, this approach yields an open-loop control, which lacks adaptability to disturbances and may fail when the system is stochastic.
% % \at{Talk about existing approaches to solve this problem. Talk about their limitations to extend to stochastic setting.}
% \begin{assumption}\label{asum-sol}
%     Problem \ref{prob:p2} has at least one solution.
% \end{assumption}
% This assumption is made to facilitate the proposed numerical algorithm.





% \at{Discuss validity of the assumptions. Add a subsection, presenting the idea of starting $\tilde x_0$ starting from a small Gaussian distribution around $y$ to ensure the assumptions are valid. Add a subsection, deriving the analytical formula for control law for linear Gaussian setting. Derive an error bound for the case where you start from a Gaussian distribution around $y$ instead of $y$. Use the result to mitivate the next subsection: improving the algorithm by adding a deterministic control.}

% We consider the time-reversal SDE
% % \begin{align}\label{eq:alg-tr-sde}
% %     \ud \tilde X_t = &-f(T-t,\tilde X_t)\ud t -\tilde g(T-t,\tilde X_t)\ud t \nonumber\\ 
% %     &- g(T-t,\tilde X_t) ({u_{T-t}} \ud t + \epsilon \ud W_t), \tilde X_0=y
% % \end{align} 
% where $u_t$ is a deterministic control input that we will specify later. $\tilde g : [0,T] \times \R^n \rightarrow \R^n$, and the $i$th component of $\tilde g$ is defined as 
% \begin{equation}
%     \tilde g_i(t,x) \mathrel{:=} \sum_{jk} g_{j,k}(t,x) \partial_{x_j}g_{i,k}(t,x)
% \end{equation}
% Let $\tilde p_0$ denote the probability density of the terminal time system state $\tilde X_T$ generated by this SDE.  The corresponding forward process admits the formulation
% \begin{align}
%     \ud X_t = f(t,X_t)\ud t + g(t,X_t) ((&{u_t}+\tilde{k}(t,X_t)) \ud t + \epsilon \ud W_t), \nonumber
%     \\ &X_0\sim \tilde{p}_0, \quad X_T = y
% \end{align}
% where $\tilde k:[0,T] \times \R^n \rightarrow \R^n$, and the $i$th component of $\tilde k$ is defined as 
% \begin{equation}
%     \tilde k_i(t,x) \mathrel{:=} \sum_j \left[ \frac{1}{p_t(x)}g_{j,i}(t,x) \partial_{x_j} p_t(x) + \partial_{x_j}g_{j,i}(t,x)\right]
% \end{equation}
% Based on our mathematical derivation in section \ref{sec:3a}, we obtain the equation:
% \begin{equation}
%     g(t,X_t) \tilde{k}(t,X_t) +\tilde g(t,X_t) = \frac{\epsilon^2}{p_t(x)} \partial_x(G(t,x)p_t(x))
% \end{equation}
%  $\tilde{k}(t,X_t)$ is the solution to minimization problem 
% \begin{equation}
%     % \min_{\tilde k(\cdot)} \mathbb{E} \left[ \frac{1}{2}\|g(t,X_t)\tilde k(t,X_t)\|^2 + \epsilon^2 \trace(G(t,X_t) \partial_x(g(t,X_t)\tilde k(t,X_t))\right]
%     \min_{\tilde k(\cdot)} \Expect \left[ \frac{1}{2}\|g\tilde k+\tilde g\|^2 + \epsilon^2 \trace(G\partial_x(g \tilde k +\tilde g))\right]
% \end{equation}
% \begin{proposition}\label{prop:sol2p1}
%     The control input $U_t = u_t + \tilde k(t,X_t)$ can bring any initial state $\tilde x$, s.t. $\tilde p_0 (X_0=\tilde x) >0$, to target state $y$ which means $U_t$ is the solution to Problem \ref{prob:p1}.
% \end{proposition}
% \begin{proof}
%     Consider a probability kernel
% \end{proof}
% Following the Proposition \ref{prop:sol2p1}, we can summarize our algorithm.
% \begin{algorithm}[bth]
% \caption{Time-reversal Algorithm} 
% \begin{algorithmic}[1]
% \STATE \textbf{Input:}  sample size $N$, step-size $\Delta t$, function class $\Psi$.   
% \STATE  $\{\tilde X_T^i\}_{i=1}^N = y$

% \FOR{$t\in \{T,T-\Delta t, \ldots,\Delta t\}$}
% \STATE $\{\Delta \tilde W^i_t\}_{i=1}^N\sim N(0,\Delta t I_n)$
% 	\STATE $\tilde X^{i}_{t-\Delta t} \!=\! \tilde X^i_{t}-(f(t,\tilde X^{i}_{t})-\epsilon^2\tilde g(t,\tilde X^{i}_{t}))\Delta t  - \epsilon g(t,\tilde X^{i}_{t}) \Delta \tilde W^i_t$
% \ENDFOR
% \STATE $k^\star(t,\cdot)\!=\! \argmin_{k\in \Psi} \frac{1}{N}\sum_{i=1}^N \big[\frac{1}{2}\|g(t,\tilde X^{i}_{t})k(t,\tilde X^i_{t}) \!+\! \tilde g(t,\tilde X^i_{t})\|^2 \!+\! \trace(G(t, \tilde X^i_{t})\partial_x (g(t,\tilde X^{i}_{t})k(t,\tilde X^i_{t})\!+\!\tilde g(t,\tilde X^i_{t})))\big]$

%  \STATE \textbf{Output:}  $\{k^\star(t,\cdot)\}_{t\in \{0,\Delta t,\ldots,T\} }$
% \end{algorithmic}
% \label{alg:naive}
% \end{algorithm}
% \begin{algorithm}[bth]
% \caption{Time-reversal Algorithm} 
% \begin{algorithmic}[1]
% \STATE \textbf{Input:}  sample size $N$, step-size $\Delta t$, control input $u_t$, function class $\Psi$.   
% \STATE  $\{\tilde X_T^i\}_{i=1}^N = y$

% \FOR{$t\in \{T,T-\Delta t, \ldots,\Delta t\}$}
% \STATE $\{\Delta \tilde W^i_t\}_{i=1}^N\sim N(0,\Delta t I_n)$
% 	\STATE $\tilde X^{i}_{t-\Delta t} \!=\! \tilde X^i_{t}-(f(t,\tilde X^{i}_{t})+\tilde g(t,\tilde X^{i}_{t})+g(t,\tilde X^{i}_{t})u_t)\Delta t  - g(t,\tilde X^{i}_{t})) \epsilon \Delta \tilde W^i_t$
% \ENDFOR
% \STATE $k(t,\cdot)\!=\! \argmin_{k\in \Psi} \frac{1}{N}\sum_{i=1}^N \big[\frac{1}{2}\|g(t,\tilde X^{i}_{t})k(t,\tilde X^i_{t}) + \tilde g(t,\tilde X^i_{t})\|^2 \!+\! \epsilon^2 \trace(G(t, \tilde X^i_{t})\partial_x (g(t,\tilde X^{i}_{t})k(t,\tilde X^i_{t})+\tilde g(t,\tilde X^i_{t})))\big]$

%  \STATE \textbf{Output:}  $\{k(t,\cdot)\}_{t\in \{0,\Delta t,\ldots,T\} }$
% \end{algorithmic}
% \label{alg:TR}
% \end{algorithm}
% \subsection{Improving Algorithm}\label{sec:ip-alg}
% The proposed algorithm, in its current form, has limited efficiency in obtaining a feedback control law due to the unforced time-reversal dynamics in \eqref{eq:alg-tr-sde}. To improve exploration efficiency, a suitable deterministic control input $u_t$ can be incorporated into the time-reversal sampling procedure. The modified time-reversal process admits the SDE 
% \begin{align}
%     \ud \tilde X_t = \tilde f(T\!-\!t,\tilde X_t)\ud t +   g(T\!-\!t,\tilde X_t)(-&u_{T\!-\!t}\ud t+ \epsilon\ud \tilde W_t),\nonumber\\&\tilde X_0=y
% \end{align}
% The corresponding forward process $X$ admits the SDE
% \begin{align}
%     \ud X_t = (-\!\tilde f(t,\!X_t)\!+\!g(t,\!X_t)u_t\!+\!&b(t,X_t)) \ud t \!+\! \epsilon g(t,\!X_t) \ud W_t, \nonumber\\&X_0\sim \tilde{p}_T, ~~X_T = y 
% \end{align}
% Following the same procedure in \ref{sec:naive-alg}, we obtain the equation
% \begin{equation}
%     f+gk = \epsilon^2 g k^\star + \epsilon^2 \tilde g - \tilde f + gu
% \end{equation}
% To satisfy the equation, we only need to design the new control law as 
% \begin{equation}
%     k(t,X_t) = u_t + \epsilon^2 k^\star(t,X_t)
% \end{equation}
% Moreover, due to time discretization and estimation errors in the score function, the feedback control law learned from the algorithm may lead to overshoot at the terminal time. To reduce terminal time errors while maintaining the same time step $\Delta t$, the time-reversal process can be initialized from a small neighborhood around $y$. Specifically, instead of setting $p_T(x) = \delta(x-y)$, it can be modified to $p_T(x) = \mathcal{N}(y, \sigma^2I)$. Based on this discussion, we present our improved algorithm
% % \begin{algorithm}[bth]
% % \caption{Improved Time-reversal Algorithm} 
% % \begin{algorithmic}[1]
% % \STATE \textbf{Input:}  sample size $N$, step-size $\Delta t$, control input $u_t$, function class $\Psi$.   
% % \STATE  $\{\tilde X_T^i\}_{i=1}^N \sim \mathcal{N}(y, \sigma^2I)$

% % \FOR{$t\in \{T,T-\Delta t, \ldots,\Delta t\}$}
% % \STATE $\{\Delta \tilde W^i_t\}_{i=1}^N\sim N(0,\Delta t I_n)$
% % 	\STATE $\tilde X^{i}_{t-\Delta t} \!=\! \tilde X^i_{t}\!-\!(f(t,\tilde X^{i}_{t})\!-\!\epsilon^2\tilde g(t,\tilde X^{i}_{t})+g(t,\tilde X^{i}_{t})u_t)\Delta t  - \epsilon g(t,\tilde X^{i}_{t})  \Delta \tilde W^i_t$
% % \ENDFOR
% % \STATE $k^\star(t,\cdot)\!=\! \argmin_{k\in \Psi} \frac{1}{N}\sum_{i=1}^N \big[\frac{1}{2}\|g(t,\tilde X^{i}_{t})k(t,\tilde X^i_{t}) \!+\! \tilde g(t,\tilde X^i_{t})\|^2 \!+\!  \trace(G(t, \tilde X^i_{t})\partial_x (g(t,\tilde X^{i}_{t})k(t,\tilde X^i_{t})\!+\!\tilde g(t,\tilde X^i_{t})))\big]$

% %  \STATE \textbf{Output:}  $\{k^\star(t,\cdot)\}_{t\in \{0,\Delta t,\ldots,T\} }$
% % \end{algorithmic}
% % \label{alg:improved}
% % \end{algorithm}
% \begin{remark}
%     A suitable candidate $u_t$ can be obtained by solving Problem \ref{prob:p2}
% \end{remark}
% \begin{proposition}\label{prop:small-noise}
%     The output $k^\star$ from Algorithm \ref{alg:naive} and \ref{alg:improved} remains valid for the same system dynamics with a different noise scalar $\epsilon^\prime$ satisfying $0\leq \epsilon^\prime \leq \epsilon$.
% \end{proposition}
% \begin{proof}
    
% \end{proof}
% \begin{remark}
%     Based on Proposition \ref{prop:sol2p1}, any control input $u_t$ that generates the $\tilde p_0$, s.t. $\tilde p_0(X_0 = x) >0$, can be applied to this algorithm. However, to enhance the efficiency of the control law, a suitable choice of $u_t$ can be obtained by solving Problem \ref{prob:p2}.
% \end{remark}
% \begin{remark}
%     Due to time discretization and estimation errors in the score function, the feedback control law learned from the algorithm may result in overshoot at the terminal time. To reduce terminal time error, the backward process can be initialized from a small neighborhood around $y$. Specifically, instead of setting $p_T(x) = \delta(x-y)$, it can be modified to $p_T(x) = \mathcal{N}(y, \sigma^2I)$.
% \end{remark}


\section{Numerical Results}\label{sec:results}
\subsection{Numerical Algorithm}
\label{sec:Algorithm}
% \begin{itemize}
%     % \item We use our proposed methodology to develop the numerical algorithm
%     % \item We simulate N random realizations of $Z_t$ SDE () In order to do so we use Euler-Maruyama approach with the step-size $\Delta t$
%     % \item In order to find the control law, we use optimization problem () where we use the decomposition () The resulting optimization problem is ...
%     % \item In order to solve this, we paramarize k with neural network, and apply ADAM. The batch size is ... uniform sample t
%     % \item the details of algorithm appears in table () 
%     % \item to illustrate the performance of the algorithm we apply in ....
% \end{itemize}
In this section, we introduce our proposed numerical algorithm which is based on the methodology described in Section~\ref{sec:methodology}. The algorithm starts with simulating $N$ random realizations $\{Z^i_t\}_{i=1}^N$ of the process \eqref{eq:Z-alg}, with the flexibility of considering an additional  deterministic control input $u_t$. The simulation is carried out using the Euler-Maruyama discretization method with the step-size $\Delta t$. In order to find the control law~\eqref{eq:fbcontrol}, we use the decomposition  of the score function \eqref{eq:score-decom} and modify the score function optimization problem \eqref{eq:score-opt} according to $\min_k\,J(gk+\mathfrak g)$. 
% \begin{align*}
%     \min_{k\in\Psi} \Expect  \Big[&\frac{1}{2} \|g(Z_t)k(t,Z_t) \!+\! \mathfrak g(Z_t)\|^2 \\
%     &+ \sum_{j,l}^nG_{j,l}(Z_t) \frac{\partial (g(Z_t)k(t,Z_t)+\mathfrak g(Z_t))_j}{\partial x_l}(t,Z_t) \Big ]
% \end{align*}
To solve this optimization problem, we parameterize $k$ with a neural network with a 3-block ResNet architecture where each block consists of 2 linear layers of width 32 and an exponential linear unit (ELU)-type activation function. We use ADAM optimizer to find the parameters of the neural network.
%, with an initial learning rate $3 \times 10^{-4}$. We use StepLR as the learning rate scheduler with a step size of $500$ and decay factor $\gamma=0.95$. 
The batch is generated by uniformly sampling $K_1=\lfloor\frac{T}{10\Delta t}\rfloor + 1$ time instants $\{t_1,t_2,\ldots,t_{K_1}\}$  from $[0,T]$, and $K_2=32$ random samples of the $N$ trajectories $\{Z^i_{t_1},\ldots,Z^i_{t_{K_1}}\}_{i=1}^N$. The details of the algorithm appear in Algorithm \ref{alg:improved}.
%To illustrate the performance of our algorithm, we apply it to three different tasks: Brownian Bridge, Inverted Pendulum, and Avoid Obstacle. 
The code for reproducing the results is available online~\footnote{\url{https://github.com/YuhangMeiUW/P2P}}.

The deterministic control input $u_t$ is designed   to approximately bring $Z_T$ to the vicinity of  $x_0$. For example, this control may be obtained by the application of trajectory optimization techniques to the deterministic version of the model~\eqref{eq:Z-alg}, when $\epsilon =0$.   In addition to decreasing the bias error, the control input $u_t$ serves as an ``importance sampling" mechanism that guides $Z_t$ to be sampled in areas of the domain where the control law is more relevant to the initial condition $x_0$, thus increasing sampling efficiency. 


% To illustrate our algorithm, we apply it to three different tasks: Brownian Bridge, Inverted Pendulum, and Avoid Obstacle. We select $\Psi$ to be the class of neural networks with  In training stage, we use ADAM optimizer, with initial learning rate $10^{-4}$ and batch size 32. In prediction stage, we use the Euler-Maruyama method. 
\subsection{Two-dimensional Brownian Bridge}\label{sec:bb-exp}
We consider a 2-dimensional system governed by the SDE
\begin{equation}\label{eq:2nd-bb}
    \ud X_t = U_t \ud t + \epsilon \ud W_t, ~X_0=x_0,
\end{equation}
and let $\epsilon=0.3$, $x_0 = (0,0)^\top$, $T=1$, and $x_f = (2,2)^\top$. This is a linear Gaussian model with $A=0$ and $B=I$. We employ the deterministic control $u_t=x_0-x_f$ which brings $Z_0=x_f$ to $Z_1=x_0$ in the deterministic setting. In this case, the resulting control law for $X_t$ takes the form 
\begin{align}
    % U_t &= u_t - \epsilon^2 B^\top Q_{1-t}^{-1}(X_t - m_{1-t}^u) \nonumber \\
    U_t=  x_f-x_0 - \frac{\epsilon^2(X_t-(1-t)x_0-tx_f)}{\epsilon^2(1-t) + \sigma^2}.\label{eq:analy-sol-u}
\end{align}
% \begin{equation}\label{eq:analy-sol}
%     U_t  = -\frac{\epsilon^2(X_t-x_f)}{\epsilon^2(1-t) + \sigma^2}
% \end{equation}
% If we follow the procedure in Section \ref{sec:improve-alg}, the analytic solution with an open-loop control $u_t = x_f-x_0$ is
% \begin{align}
%     % U_t &= u_t - \epsilon^2 B^\top Q_{1-t}^{-1}(X_t - m_{1-t}^u) \nonumber \\
%     U_t= - \frac{\epsilon^2(X_t-x_f) -\sigma^2x_f}{\epsilon^2(1-t) + \sigma^2}\label{eq:analy-sol-u}
% \end{align}

We apply Algorithm \ref{alg:improved} with $u_t = x_0-x_f$, $N=1000$, $\sigma=0$, and $\Delta t = 0.004$. The resulting  trajectories $\{Z^i_t\}_{i=1}^N$ and $\{X^i_t\}_{i=1}^N$, along with the control inputs $\{U^i_t\}_{i=1}^N$,  are shown in Fig. \ref{fig:bb-X-compare}. 
% For comparison, we show the trajectories that only incorporate the open-loop input $\tilde u_t = x_f-x_0$.  
% Want to find a feedback control law $U_t = k(t,X_t)$ such that $X_{1} = x_f$. This problem is known as the 2nd order Brownian Bridge.  
% We apply Algorithm \ref{alg:improved} with $u_t = tx_0+(1-t)x_f$. In this experiment, we set the initial state and the target state as $x_0 = (0,0)^\top$ and $x_f = (2,2)^\top$. 
% The comparison of our control law and open loop control 
The result demonstrates that the control law obtained from Algorithm \ref{alg:improved} successfully steers all trajectories $X^i_t$ from $x_0$ to desired target $x_f$.
% \begin{figure}[tb]
%     \centering
%     \includegraphics[width=0.7\hsize,trim={8pt 0pt 30pt 40pt},clip]{figs/NNu_Openloop_sigma00_epsilon03_N1000_T251.pdf}
%     \caption{Numerical result for Brownian Bridge. $x_0=(0,0)^\top$ and $x_f=(2,2)^\top$ The time horizon $T = 1$, time step-size $\Delta t=0.004$, sample size $N = 1000$, and noise scalar $\epsilon=0.3$. The trajectories generated by our algorithm are plotted, with those under open-loop control shown for comparison.}
%     \label{fig:bb-X-compare}
% \end{figure}
\begin{figure*}[tb]
	\centering
    \begin{subfigure}{0.32\textwidth}\centering
        \includegraphics[width=\hsize,trim={0pt 0pt 12pt 41pt},clip]{figs/X_Z1_traj_new.pdf}
        \caption{}
        \label{fig:xz1}
    \end{subfigure}
    \begin{subfigure}{0.32\textwidth}
        \includegraphics[width=\hsize,trim={0pt -5pt 12pt 41 pt},clip]{figs/X_Z2_traj_new.pdf}
        \caption{}
        \label{fig:xz2}
    \end{subfigure}
    \begin{subfigure}{0.32\textwidth}
        \includegraphics[width=\hsize,trim={0pt 0pt 12pt 41 pt},clip]{figs/u_traj_new.pdf}
        \caption{}
        \label{fig:u}
    \end{subfigure}
	\caption{Numerical result for the application Algorithm~\ref{alg:improved} to the two-dimensional Brownian bridge example of Section~\ref{sec:bb-exp}: (a) First component of $X_t$ and $\tilde Z_t$ (b) Second component of $X_t$ and $\tilde Z_t$ (c) Control input $U_t$.
    %$x_0=(0,0)^\top$ and $x_f=(2,2)^\top$. The time horizon $T = 1$, time step-size $\Delta t=0.004$, sample size $N = 1000$, noise scalar $\epsilon=0.3$, and $\sigma = 0$.
    }
    \label{fig:bb-X-compare}
    %Time horizon $T=1$, sample size $N=1000$, and noise scalar $\epsilon=0.3$}
\end{figure*}


In order to quantify the performance of the control law, we introduce the following mean-squared-error (MSE) criteria 
\begin{equation}\label{eq:mse}
    MSE = \frac{1}{N}\sum_{i=1}^N \|X_T^i - x_f\|_2^2.
\end{equation}
% where $\{X^i_t\}_{i=1}^N$ are the sample trajectories of $X_t$. 
% We evaluate the MSE for the following methods: 
% \begin{itemize}
%     \item NN+$u_t$: Use a neural network to approximate $k^\star$ in Algorithm \ref{alg:improved} with $u_t=x_f-x_0$
%     \item NN: Use a neural network to approximate $k^\star$ in Algorithm \ref{alg:improved} with $u_t=0$
%     \item Sol: Use the analytic solution \eqref{eq:analy-sol}
%     \item Sol+$u_t$: Use the analytic solution \eqref{eq:analy-sol-u}
% \end{itemize}
% \at{Describe the figures one by one.}
We investigate the influence of the time step $\Delta t$ and standard deviation $\sigma$ on the MSE criteria. 
% In order to do so, we evaluate the MSE \eqref{eq:mse} averaged over 5 experiments for $\Delta t \in \{0.002, 0.004, 0.01, 0.02, 0.05, 0.1\}$ and $\sigma \in \{0.0, 0.02, 0.04, 0.05, 0.1, 0.2, 0.3\}$. The open-loop control $u_t=x_f-x_0$ is used as a baseline for comparison. 
The result for varying time step-size is presented in Fig \ref{fig:mse_dt}, where we fix $\sigma=0$. We also show the MSE corresponding to implementing the exact form of the control law~\eqref{eq:analy-sol-u} and the open-loop control $U_t=x_f-x_0$ as baselines. It is observed that the algorithm performs almost as good as the exact solution, and the MSE decreases as $\Delta t \to 0$. The result for varying $\sigma$ is presented in Fig \ref{fig:mse_sigma} with fixed $\Delta t = 0.004$, where for comparison, the case without using the deterministic input is also included. It is observed that including the deterministic input significantly decreases the MSE when $\sigma>0$, justifying the remark~\ref{rem:improve-alg}. Moreover, $\sigma$ acts as a regularizer in the optimization procedure and decreases the difference between Algorithm 1 and the exact solution. Although increasing $\sigma$ increases MSE, it serves to avoid singularity of the control,  which is shown by computing  the averaged control energy 
\begin{equation}\label{eq:unorm}
    U_{norm} = \frac{1}{N}\sum_{i=1}^N\int_0^T \|U_t^i\|^2 \ud t
\end{equation}
as a function of $\sigma$, in Fig \ref{fig:unorm}. The MSE and  $U_{norm}$ are both averaged over $5$ independent experiments, where the shaded region represents the range from the minimum to the maximum across experiments.
\begin{figure*}[tb]
	\centering
    
    \begin{subfigure}{0.32\textwidth}\centering
        \includegraphics[width=\hsize,trim={0pt 0pt 12pt 41pt},clip]{figs/MSE_1divdt_epsilon03_N1000_sigma00.pdf}
        \caption{}
        \label{fig:mse_dt}
    \end{subfigure}
    \begin{subfigure}{0.32\textwidth}
        \includegraphics[width=\hsize,trim={0pt -5pt 12pt 41 pt},clip]{figs/MSE_sigma_epsilon03_N1000_T251.pdf}
        \caption{}
        \label{fig:mse_sigma}
    \end{subfigure}
    \begin{subfigure}{0.32\textwidth}
        \includegraphics[width=\hsize,trim={0pt 0pt 12pt 41 pt},clip]{figs/unorm_sigma_epsilon03_N1000_T251.pdf}
        \caption{}
        \label{fig:unorm}
    \end{subfigure}
	\caption{Numerical error analysis for the application Algorithm~\ref{alg:improved} to the two-dimensional Brownian bridge example of Section~\ref{sec:bb-exp}:  (a) Influence of time step-size on MSE~\eqref{eq:mse}, with $\sigma=0$ (b) Influence of $\sigma$ on MSE, with $\Delta t = 0.004$ (c) Influence of $\sigma$ on $U_{norm}$~\eqref{eq:unorm}, with $\Delta t = 0.004$. The results compare (i) Algorithm \ref{alg:improved}; (ii) the exact solution~\eqref{eq:analy-sol-u}; and  implementing the open loop control $U_t=x_f-x_0$. Panel (b) also includes the case where the deterministic input $u_t=0$ in both exact solution and Algorithm~\ref{alg:improved}.}
    %Time horizon $T=1$, sample size $N=1000$, and noise scalar $\epsilon=0.3$}
\end{figure*}
% \begin{figure*}[tb]
% 	\centering
% 	\includegraphics[width=0.32\hsize,trim={0pt 0pt 12pt 41 pt},clip]{figs/MSE_dt_epsilon03_N1000_sigma00.pdf} \hfill 
% 	\includegraphics[width=0.32\hsize,trim={0pt 0pt 12pt 41 pt},clip]{figs/MSE_sigma_epsilon03_N1000_T251.pdf} \hfill 
% 	\includegraphics[width=0.32\hsize,trim={0pt 0pt 12pt 41 pt},clip]{figs/unorm_sigma_epsilon03_N1000_T251.pdf}\\
% 	\hspace{20pt}(a) \hspace{160pt} (b)  \hspace{160pt}  (c) 
% 	\caption{Numerical result for .}
% 	\label{fig:dt-sig-influence}
% \end{figure*}
% \at{Include problem parameters: $\epsilon$, $t_f$, \ldots. Include the analytical expressions for the control law. Look at the flow matching paper for how to present the results. Define MSE error and include error result as a function of $N$, $\Delta t$. Describe the different ways you compute the control law. Compare the result to the case where the deterministic control is not included. Include comparison for different values of $\sigma$.}
\subsection{Inverted Pendulum}\label{sec:IP-exp}
We consider the  stochastic pendulum dynamics with 
\begin{align*}
    f(x) =  \begin{bmatrix}
        x(2) \\
        \sin(x(1)) - 0.01x(2)
    \end{bmatrix},\quad g(x) = \begin{bmatrix}
        0\\1
    \end{bmatrix},
\end{align*}
% system $X_t = \begin{bmatrix}
%     X_1(t) \\
%     X_2(t)
% \end{bmatrix}$, where $X_1(t)$ denotes angle and $X_2(t)$ denotes angle velocity, governed by SDE
% \begin{align}
%     \ud X_t = \begin{bmatrix}
%         X_t(2) \\
%         \sin(X_t(1)) - 0.01X_t(2)
%     \end{bmatrix} \ud t + &\begin{bmatrix}
%         0\\1
%     \end{bmatrix} (U_t \ud t +\epsilon \ud W_t), \nonumber \\
%     &X_0 = \begin{bmatrix}
%         \pm \pi \\0
%     \end{bmatrix}
% \end{align}
% where 
$\epsilon=0.3$, $x_0=[\pi,0]^\top$, $x_f = [0,0]^\top$, and $T=5$. The first component of the state is the angle where the angle equal to $0$ denotes the upward position of the pendulum and $\pi$ denote the downward. The goal is to bring the pendulum from the downward position to the upward position. We apply Algorithm~\ref{alg:improved} with $N=1000$, $\sigma=0$, $u_t=0$, and $\Delta t =0.004$. The resulting trajectories and control inputs are shown in Fig. \ref{fig:ip-traj}, where the successful steering of the pendulum to the upward position is demonstrated. 

% Moreover, we consider an extension of the pendulum dynamics with
% \begin{align*}
%     f(x) = \begin{bmatrix}
%         x(2) \\
%         \sin(x(1)) - 0.01x(2) + x(3)\\
%         0
%     \end{bmatrix}, \quad g(x) = \begin{bmatrix}
%         0\\0\\1
%     \end{bmatrix}. 
% \end{align*}
% In this model, the torque that is applied to the pendulum is modeled as the third component of the state $x(3)$, and the control input $U_t$ affects  the change in the torque, as opposed to the torque itself in the original model. We consider the same parameters as before, with the difference of initializing $Z_0(3)\sim \mathcal{N}(\cdot; 0,\sigma^2)$ with $\sigma = 0.05$. 
% %The first and second components of the state are defined in the same way as the 2D inverted pendulum. The third component of the state is the control input for the 2D inverted pendulum. 
% The resulting trajectories from the application of Algorithm \ref{alg:improved} are shown in Fig. \ref{fig:3d-ip}. It is observed that the resulting torque values $X_t(3)$ that brings pendulum to the upward position are smaller compared to torques  from the original model. 
% For comparison, we implemented the linear-quadratic regulator on the  model linearized around $x=[0,0]^\top$, fine-tuning its parameters to given the least MSE. The comparison is reported in Table~\ref{}. 
% The initial state is the bottom and the target state is top. We use trajectory optimization method to find an open loop control which steers the inverted pendulum from $x_0$ to $x_f$ in the deterministic setting.

% \begin{figure}[htp]
%     \centering
%     \includegraphics[width=0.9\linewidth]{figs/pendulum_backward.pdf}
%     \caption{Trajectories used to train. The time horizon $T = 8$, time step-size $\Delta t=0.004$, and sample size $N = 2000$, noise scalar $\epsilon=0.3$}
%     \label{fig:pendulum-back}
% \end{figure}

% \begin{figure}[htp]
%     \centering
%     \includegraphics[width=0.7\linewidth]{figs/IP_NNu_Openloop_sigma00_epsilon03_N1000_T251.pdf}
%     \caption{Numerical result for Inverted Pendulum. Compare to open loop input}
%     \label{fig:Inverted-pendulum}
% \end{figure}
\begin{figure*}[tb]
	\centering
    \begin{subfigure}{0.32\textwidth}\centering
        \includegraphics[width=\hsize,trim={0pt 0pt 12pt 41pt},clip]{figs/IP_X_Z1_traj_new.pdf}
        \caption{}
        \label{fig:ip-xz1}
    \end{subfigure}
    \begin{subfigure}{0.32\textwidth}
        \includegraphics[width=\hsize,trim={0pt -5pt 12pt 41 pt},clip]{figs/IP_X_Z2_traj_new.pdf}
        \caption{}
        \label{fig:ip-xz2}
    \end{subfigure}
    \begin{subfigure}{0.32\textwidth}
        \includegraphics[width=\hsize,trim={0pt 0pt 12pt 41 pt},clip]{figs/IP_u_traj_new.pdf}
        \caption{}
        \label{fig:ip-u}
    \end{subfigure}
	\caption{Numerical result for the application Algorithm~\ref{alg:improved} to the inverted pendulum example of Section~\ref{sec:IP-exp}: (a) First component of $X_t$ and $\tilde Z_t$ ; (b) Second component of $X_t$ and $\tilde Z_t$; (c) Control input $U_t$. The first component of the initial state  is equal to $\pm\pi$ and represents the downward position of the pendulum, while the first component of the terminal state is equal to $0$, representing the upward position. 
    %$x_0=[\pi,0]^\top$ and $x_f=[0,0]^\top$. The time horizon $T = 5$, time step-size $\Delta t=0.004$, sample size $N = 1000$, and noise scalar $\epsilon=0.3$.
    }
    \label{fig:ip-traj}
    %Time horizon $T=1$, sample size $N=1000$, and noise scalar $\epsilon=0.3$}
\end{figure*}

% \begin{figure*}[tb]
% 	\centering
%     \begin{subfigure}{0.32\textwidth}\centering
%         \includegraphics[width=\hsize,trim={0pt 0pt 12pt 41pt},clip]{figs/3dIP_X_Z1_traj_new.pdf}
%         \caption{}
%         \label{fig:3ip-xz1}
%     \end{subfigure}
%     \begin{subfigure}{0.32\textwidth}
%         \includegraphics[width=\hsize,trim={0pt -5pt 12pt 41 pt},clip]{figs/3dIP_X_Z2_traj_new.pdf}
%         \caption{}
%         \label{fig:3ip-xz2}
%     \end{subfigure}
%     \begin{subfigure}{0.32\textwidth}
%         \includegraphics[width=\hsize,trim={0pt 0pt 12pt 41 pt},clip]{figs/3dIP_X_Z3_traj_new.pdf}
%         \caption{}
%         \label{fig:3ip-xz3}
%     \end{subfigure}
% 	\caption{Numerical result for the application Algorithm~\ref{alg:improved} to the extend version of the inverted pendulum example in Section~\ref{sec:IP-exp}:  \ref{sec:IP-exp}: (a) First component of $X_t$ and $\tilde Z_t$ (b) Second component of $X_t$ and $\tilde Z_t$ (c) Third component of $X_t$ and $\tilde Z_t$. 
%     %$x_0=[\pi,0, \mathcal N(0,\sigma^2I)]^\top$ and $x_f=[0,0,\mathcal N(0,\sigma^2I)]^\top$. The time horizon $T = 5$, time step-size $\Delta t=0.004$, sample size $N = 1000$, noise scalar $\epsilon=0.3$, and $\sigma=0.05$.
%     }
%     \label{fig:3d-ip}
%     %Time horizon $T=1$, sample size $N=1000$, and noise scalar $\epsilon=0.3$}
% \end{figure*}

% \begin{figure}
%     \centering
%     \includegraphics[width=0.9\linewidth]{figs/inverted_pendulum_nonoise.pdf}
%     \caption{Numerical result for same feedback controller for deterministic Inverted Pendulum, i.e. $\epsilon=0$}
%     \label{fig:IP-nonoise}
% \end{figure}

% \subsection{Avoid Obstacle}

% Consider the same system defined in \eqref{eq:2nd-bb}, we want to find a feedback control law such that $X_1 = y$ and avoid obstacle in the trajectory. We simulate the time-reversal process starting from $y$ and eliminate trajectories that intersect with the obstacle. To improve the efficiency, we also remove the trajectories that are closed to the obstacle.

% \begin{figure}[htp]
%     \centering
%     \includegraphics[width=0.9\linewidth]{figs/elimited_trajectories.pdf}
%     \caption{Trajectories used to train a controller. Obstacle is a circle center at $(1,1)$ with radius 0.2. We use a bigger radius 0.6 to eliminate the trajectories to improve efficiency}
%     \label{fig:enter-label}
% \end{figure}

% The feedback control law apply to 15000 samples, and the success rate is 0.9979. In the deterministic setting the same controller success rate is 1.0.

% \begin{figure}[htp]
%     \centering
%     \includegraphics[width=0.9\linewidth]{figs/avoid_obstacle_deeper.pdf}
%     \caption{Numerical result for Avoid Obstacle. The time horizon $T = 1$, time step-size $\Delta t=0.001$, and sample size $N = 15000$, noise scalar $\epsilon=0.5$.}
%     \label{fig:avoid-obstacle}
% \end{figure}

% \begin{figure}
%     \centering
%     \includegraphics[width=0.9\linewidth]{figs/avoid_obstacle_nonoise.pdf}
%     \caption{Numerical result for Avoid Obstacle in deterministic setting}
%     \label{fig:avoid-nonoise}
% \end{figure}



{
\section{Conclusion}

This paper introduces a novel approach for steering nonlinear stochastic control-affine systems to a desired target state within a finite time horizon, leveraging time-reversal theory of diffusions. By constructing feedback control laws based on the score function associated with the reversed dynamics, the proposed method ensures finite-time convergence to the target state. Unlike traditional Schr\"{o}dinger bridge methods or stochastic optimal control formulations, our approach is computationally efficient and applicable to both linear and nonlinear stochastic systems without relying on optimality relative to the initial condition.

An extension of the theory to address practical challenges related to inevitable singularities in the control law near the terminal time is also presented through relaxation of the almost-sure constraint to a distribution constraint and explicit error bounds of this approach are provided for the linear Gaussian case. Numerical experiments demonstrate the effectiveness of the method across benchmark examples, including a Brownian bridge and inverted pendulum dynamics. 

Future research includes the extension of the theoretical results for assigning terminal distribution constraints in the general nonlinear system setting, the improving score function approximation techniques,  exploring the optimality properties of the proposed control law through its connection to Doob's $h$-transform, and numerical verification of our method on real-world control applications in Aerospace~\cite{exarchos2019optimal,lew2020chance}.





}

% \section{Error Analysis}\label{sec:err-analy}
% \input{error_analysis}

    \bibliographystyle{IEEEtran}
    \bibliography{references}

     % \appendix
     % \input{appendix.tex}
\end{document}

